\documentclass[10pt,a4paper]{amsart}
\usepackage{amsmath,amssymb,amsthm,graphicx}
\numberwithin{equation}{section}

\newtheorem{Theorem}{Theorem}[section]
\newtheorem*{Theorem*}{Theorem}
\newtheorem{Corollary}[Theorem]{Corollary}
\newtheorem{Lemma}[Theorem]{Lemma}
\newtheorem{Proposition}[Theorem]{Proposition}

\theoremstyle{definition}
\newtheorem{Definition}[Theorem]{Definition}
\newtheorem{Note}[Theorem]{Note}
\newtheorem{Example}[Theorem]{Example}
\newtheorem{Remark}[Theorem]{Remark}

%%%%%%%%%%%%%%% My definitions %%%%%%%%%%%%%%%%%%%%%%%
\def\nn{\notag}
\def\Z2{\mathbb{Z}_2^2}
\def\g{\mathfrak{g}}
\def\h{\mathfrak{h}}
\def\DP#1#2{\hat{#1}\cdot \hat{#2}}
\def\PH#1#2{(-1)^{\hat{#1}\cdot \hat{#2}}}
\def\osp{\mathfrak{osp}}
\def\gl{\mathfrak{gl}}
\def\sl{\mathfrak{sl}}
\def\TL{\mathcal{L}}
\def\qq{$\Z2$-$\q(n)$}
\def\q{\mathfrak{q}}
\def\so{\mathfrak{so}}
\def\sp{\mathfrak{sp}}
\def\Zsl{$\mathbb{Z}_3^2$-$\sl(2)$}
\def\Zosp{$\Z2$-$\osp(m|2n)$}
%%%%%%%%%%%%%%%%%%%%%%%%%%%%%%%%%%%%%%%%%%%%%%%%%%%%%%%

\usepackage{bm}%
\usepackage{stmaryrd}
\usepackage{nicematrix}
\usepackage{arydshln}
\usepackage{braket}
\usepackage[bookmarks=false]{hyperref}
\usepackage{cleveref}
\crefname{originalreference}{original reference}{original references}
\Crefname{originalreference}{Original reference}{Original references}
\makeatletter
\newlabel{SEC:qn}{{4.1}{source}{}{sourcepointer-1}{}}
\newlabel{SEC:qn@cref}{{[originalreference][0][]4.1}{[1][0][]source}{}{}{}}
\makeatother
\title{A doubled color osp(m|2n) family with nonzero-degree Casimir elements}
\author{Naruhiko Aizawa, Ichi Fujii, Jambulingam Segar, Joris Van der Jeugt}
\date{doi:10.3842/SIGMA.2026.070; CC BY 4.0}
\begin{document}
\maketitle
\noindent\textit{Allowed standard prerequisites (editorial).}\begin{itemize}
\item Representations, color Lie algebras and universal enveloping algebras. Original Scheunert and color-algebra citations retained in full bibliography and preliminaries.
\end{itemize}\medskip
\ifdefined\hypertarget\hypertarget{sourcepointer-1}{}\fi
\noindent\textit{Source reference.} Original reference 4.1: algebra-geometry/sources/sigma-2026-070/native/sigma26-070.tex, source lines 620--620. Author narrative comparison or example pointer; complete construction used for this target is included.\par\medskip
\section{Preliminaries} \label{SEC:Pre}

This section presents the basic notions of color Lie algebras that will be used in this paper.
We give only the minimum necessary for the present study.
The reader may refer to \cite{scheu,scheuTen,zhang2025} for further details.

Let $\Gamma$ be an additive Abelian group. We fix $\Gamma$ and suppose that all vector spaces and algebras are defined over the field $\mathbb{C}$.
A vector space $V$ is said to be $\Gamma$-graded if it can be written as a~direct sum of subspaces indexed by elements of $\Gamma$:
\begin{equation*}
	V = \bigoplus_{\gamma \in \Gamma} V_{\gamma}.
\end{equation*}
We say that an element of $ V_{\gamma}$ has \textit{degree} $\gamma$.
A linear mapping $ f\colon V \to W $ between two $\Gamma$-graded vectors spaces $V$ and $ W $ is said to be of degree $\gamma \in \Gamma $ if
\begin{equation*}
	f(V_{\alpha}) \subset W_{\alpha+\gamma}, \qquad \forall \alpha \in \Gamma.
\end{equation*}
The vector space $ \operatorname{Hom}	(V,W)$ of all linear mappings of $ V $ into $ W $ is naturally $\Gamma$-graded.

An algebra $S$ is called $ \Gamma$-graded if its underlying vector space is $\Gamma$-graded and if it satisfies
\begin{equation*}
	S_{\alpha} S_{\beta} \subset S_{\alpha+\beta}, \qquad \alpha, \beta \in \Gamma.
\end{equation*}
To define color Lie algebras ($\Gamma$-graded $\omega$-Lie algebras), we introduce a mapping called a \textit{commutative factor} on $\Gamma$ (called commutation factor in \cite{scheu}) $ \omega \colon \Gamma \times \Gamma \to \mathbb{C}^{\ast}$, which satisfies
\begin{align}
	&\omega(\alpha,\beta)= \omega(\beta,\alpha)^{-1},
	\nn \\
	&\omega(\alpha,\beta+\gamma)= \omega(\alpha,\beta) \omega(\alpha,\gamma),
	\nn \\
	&\omega(\alpha+\beta,\gamma)= \omega(\alpha,\gamma) \omega(\beta,\gamma), \qquad
	\forall \alpha, \beta, \gamma.
	\label{omegaprop1}
\end{align}
These relations imply
\begin{align}
	&\omega(0,0)= \omega(0,\alpha) = \omega(\alpha,0) =1, \qquad \omega(\alpha,\alpha) = \pm 1,
	\nn\\
	&\omega(\alpha,-\beta)= \omega(-\alpha,\beta) = \omega(\beta,\alpha), \qquad
	\omega(-\alpha,-\beta) = \omega(\alpha,\beta), \qquad \forall \alpha, \beta.
	\label{omegaprop2}
\end{align}

\begin{Definition}
	Let \smash{$\g = \bigoplus_{\gamma\in\Gamma}\g_{\gamma}$} be a $\Gamma$-graded vector space and let $\omega$ be a commutative factor on $\Gamma.$
	Then $\g$ is called a color Lie algebra, if $\g$ admits a bilinear map
	$ \llbracket\ , \ \rrbracket\colon \g \times \g \to \g, $ which is of degree $0$, satisfying the following conditions:
	\begin{align}
		&\llbracket X,Y \rrbracket= -\omega(\alpha,\beta)\llbracket Y,X\rrbracket,
		\nonumber\\
		&\llbracket X, \llbracket Y, Z \rrbracket \rrbracket= \llbracket \llbracket X, Y \rrbracket, Z \rrbracket + \omega(\alpha,\beta) \llbracket Y, \llbracket X, Z \rrbracket \rrbracket,
		\qquad \forall X, Y, Z \in \g,
		\label{colorJacobi}
	\end{align}
	where $ \alpha = \deg(X)$, $\beta = \deg(Y)$ are degrees of the elements $ X$, $Y$, respectively.
\end{Definition}	

We call the bilinear map $ \llbracket\ , \ \rrbracket $ a color Lie bracket.
In some literature, a distinction is made between color Lie algebras and color Lie superalgebras.
The difference lies in the sign of $ \omega(\alpha,\alpha)$.
If $ \omega(\alpha,\alpha) = 1 $ for all $ \alpha \in \Gamma,$ then $\g$ is called a color Lie algebra. If there exists at least one $\alpha \in \Gamma$ such that $\omega(\alpha,\alpha) = -1$, then $\g$ is called a color Lie superalgebra. We do not make this distinction in this paper.

Any $\Gamma$-graded associative algebra $ S$ may be turned into a color Lie algebra by defining the color Lie bracket as
\begin{equation}
	\llbracket X, Y \rrbracket = XY - \omega(\alpha,\beta) YX, \qquad
	X, Y \in S, \quad \alpha = \deg(X),\quad \beta = \deg(Y)
	\label{ColorBracket}
\end{equation}
for any commutative factor $\omega.$
In the case of $ \omega(\alpha,\beta) = +1 $ or $ -1$ for all $\alpha, \beta \in \Gamma, $ the color Lie bracket is a commutator or an anticommutator.

Like for Lie algebras and Lie superalgebras, $ \operatorname{End}(V) := \operatorname{Hom}(V,V)$ on a $\Gamma$-graded vector space $V$ provides an example of color Lie algebras.
Observe that $\operatorname{End}(V)$ is an associative $\Gamma$\nobreakdash-graded algebra with the composition of endomorphisms as the multiplication.
Then, \textit{the general linear color Lie algebra} $ \gl(V,\omega)$ is defined as $\operatorname{End}(V)$ equipped with the color Lie bracket given by~\eqref{ColorBracket}.
To be more explicit, we present a basis of $\gl(V,\omega)$ in terms of matrices.
Suppose that $ \alpha < \beta < \gamma < \cdots $ is a fixed total order for $\Gamma.$
Choose an ordered basis for each graded subspace of $V$ and order the basis of $ V_{\alpha}, V_{\beta}, V_{\gamma}, \dots $ according to the total order for $\Gamma.$
This yields an ordered basis of~$V$.
Then an element $E(\alpha,\beta) \in \operatorname{Hom}(V_{\beta}, V_{\alpha})$ is represented by a block matrix of the following form:
\begin{equation}
	E(\alpha,\beta)=
	\bordermatrix{ & \alpha & \beta & \gamma & \cdots \cr
		\alpha & \bm{\cdot} & \ast & \bm{\cdot} & \ldots \cr
		\beta & \bm{\cdot} & \bm{\cdot} & \bm{\cdot} \cr
		\gamma & \bm{\cdot} & \bm{\cdot} & \bm{\cdot} & \ldots \cr
		\vdots & \vdots & \vdots & \vdots & \cr
	},
	\label{Edef}
\end{equation}
where the dot $\bm{\cdot}$ denotes a zero matrix and the asterisk $\ast$ denotes the block containing nonzero entries.
We denote the matrix unit relative to the ordered basis of $V$ by $ E(\alpha,\beta)_{ij} $ where $ 1 \leq i \leq \dim V_{\alpha}$, $1 \leq j \leq \dim V_{\beta}$. It is the matrix~\eqref{Edef} whose only nonzero entry is $1$ at the $(i,j)$-position of the asterisk block.
It is easy to see the following:
\begin{align*}
	&\deg(E(\alpha,\beta)_{ij})= \alpha - \beta,
	\qquad
	E(\alpha,\beta)_{ij} E(\gamma,\delta)_{k\ell}= \delta_{\beta\gamma} \delta_{jk} E(\alpha,\delta)_{i\ell}.
\end{align*}
It follows that the matrix units form a basis of $\gl(V,\omega)$ and we can express the defining relations of $\gl(V,\omega)$ as
\begin{equation*}
	\llbracket E(\alpha,\beta)_{ij}, E(\gamma,\delta)_{k\ell} \rrbracket =
	\delta_{\beta\gamma} \delta_{jk} E(\alpha,\delta)_{i\ell} - \omega(\alpha-\beta,\gamma-\delta) \delta_{\alpha\delta} \delta_{\ell i} E(\gamma,\beta)_{kj}.
\end{equation*}

The degree $\alpha$ subspace of $\gl(V,\omega)$ is
$ \gl(V,\omega)_{\alpha} := \sum_{\beta \in \Gamma} \operatorname{Hom}(V_{\beta}, V_{\beta+\alpha}).$
An element $ X(\alpha,\beta) \in \gl(V,\omega)_{\alpha-\beta}$ is a linear combination of $E(\alpha,\beta)_{ij}$:
\begin{equation*}
	X(\alpha,\beta) = \sum_{i=1}^{\dim V_\alpha} \sum_{j=1}^{\dim V_\beta} x_{ij} E(\alpha,\beta)_{ij}, \qquad x_{ij} \in \mathbb{C},
\end{equation*}
and an element $ X \in \gl(V,\omega)$ can be written as
\begin{equation*}
	X = \sum_{\alpha,\beta\in \Gamma} X(\alpha,\beta).
\end{equation*}
We now define a generalization of the trace.

\begin{Definition}
	The color trace $\operatorname{ctr}\colon \operatorname{End}(V) \to \mathbb{C}$ is defined by
	\begin{equation*}
		\operatorname{ctr} X = \sum_{\alpha \in \Gamma} \omega(\alpha,\alpha) \operatorname{tr} X(\alpha,\alpha),
		%\label{ctrdef}
	\end{equation*}
	where $\operatorname{tr}$ denotes the usual trace of a matrix. Clearly, the color trace is a mapping of degree zero.
\end{Definition}
%
\begin{Proposition}[\cite{scheuTen,zhang2025}] We have
	\begin{equation}
		\operatorname{ctr}(X(\alpha) X(\beta)) = \omega(\alpha,\beta) \operatorname{ctr}(X(\beta) X(\alpha)),
		\label{ctrperm}
	\end{equation}
	where $ X(\alpha)$ and $ X(\beta) $ are homogeneous elements of $\gl(V,\omega)$ of degree $\alpha$ and $\beta$, respectively.
\end{Proposition}

Obviously,
$ \sl(V,\omega) := \{ X \in \gl(V,\omega)\mid \operatorname{ctr} X = 0 \} $
is a $\Gamma$-graded subalgebra of $\gl(V,\omega). $
To introduce $\Gamma$-graded subalgebras of $\gl(V,\omega)$, we introduce a bilinear form on~$V$.
Let $ J\colon V \times V \to \mathbb{C}$ be a degree-zero and nondegenerate bilinear form on~$V$.
The condition $\deg(J) = 0 $ implies that for any $ x, y \in V$
\begin{equation*}
	J(x,y) = 0 \qquad \text{if} \ \deg(x) + \deg(y) \neq 0.
\end{equation*}
The nondegeneracy of $J$ requires that $ \dim V_{\alpha} = \dim V_{-\alpha} $ for all $\alpha \in \Gamma$.
This can be seen by representing $J$ in matrix form relative to the ordered basis of~$V$.
Suppose that $ \alpha \neq -\alpha $ and $ \alpha' = -\alpha'$ for $ \alpha, \alpha' \in \Gamma,$ and they are ordered as $ \cdots < \alpha < \cdots < \alpha' < \cdots < - \alpha < \cdots $.
Then $J$ has the following block matrix structure:
\begin{equation*}
	J =
	\bordermatrix{ & \alpha & \cdots & \alpha' & \cdots& -\alpha & \cdots \cr
		\alpha & \bm{\cdot} & \cdots & \bm{\cdot} & \cdots & \ast & \cdots \cr
		\; \vdots & \vdots & & \vdots & & \vdots & \cdots \cr
		\alpha' & \bm{\cdot} & \cdots & \ast & \cdots & \bm{\cdot} & \cdots \cr
		\; \vdots & \vdots & & \vdots & & \vdots & \cdots \cr
		-\alpha & \ast & \cdots & \bm{\cdot} & \cdots & \bm{\cdot} & \cdots \cr
		\;\vdots & \vdots & & \vdots & & \vdots &
	},
	%\label{JmatrixGeneral}
\end{equation*}
where the symbols $ \bm{\cdot}$ and $\ast$ have the same meaning as in \eqref{Edef}.
Note that the rows and columns intersecting a $\ast$ block contain only zero entries.
If $\dim V_{\alpha} \neq \dim V_{-\alpha},$ then the $(\alpha,-\alpha)$ and $(-\alpha,\alpha)$ blocks are rectangular matrices (not square).
This implies that $\det J = 0$, since one necessarily selects a zero entry when choosing one entry from each row and each column.
This shows that $J$ is degenerate.

We are now in a position to define subalgebras of $\gl(V,\omega)$ using the bilinear form $J.$

\begin{Proposition}[\cite{scheuTen}]\label{Prop:SulAlg}
	Let $ u$, $v$ be homogeneous elements of $V$ and $ J$ be a degree-zero nondegenerate bilinear form on $V.$
	Suppose that $J$ is either symmetric or skew-symmetric.
	Then the set of all homogeneous elements of $\gl(V,\omega)$ satisfying
\begin{equation}
J(Xu,v)+ \omega(\gamma,\alpha) J(u,Xv) = 0, \qquad X \in \gl(V,\omega)_{\gamma},\quad \forall u \in V_{\alpha}, \quad v \in V_{\beta}	\label{SubalgCond}
\end{equation}
	forms a $\Gamma$-graded subalgebra of $\gl(V,\omega)$.
\end{Proposition}

\begin{proof}
	It is easy to verify that if $ X \in \gl(V,\omega)_{\gamma} $ and $ Y \in \gl(V,\omega)_{\delta}$ satisfy \eqref{SubalgCond}, then $ \llbracket X, Y \rrbracket \in \gl(V,\omega)_{\gamma+\delta}$ also satisfies~\eqref{SubalgCond}.
	
	Next, consider the relation obtained by exchanging $ u $ and $ v $ in \eqref{SubalgCond}:
	\begin{equation}
		J(Xv,u) + \omega(\gamma,\beta) J(v, Xu) = 0. \label{SubalgCond2}
	\end{equation}
	A sufficient condition for the compatibility of this relation with \eqref{SubalgCond} is
	\begin{equation}
		J(u,v) = \pm \omega(\alpha,\beta) J(v,u). \label{Jsymm}
	\end{equation}
	Using this relation and the properties of the commutative factor given in \eqref{omegaprop1} and \eqref{omegaprop2}, one can see that the left hand side of \eqref{SubalgCond2} reduces to that of \eqref{SubalgCond}.
	Recall that $ J$ is nonvanishing if and only if $ \alpha + \beta = 0.$
	In this case, \eqref{Jsymm} yields
	\begin{equation*}
		J(u,v) = \pm \omega(\alpha,-\alpha) J(v,u) \stackrel{\eqref{omegaprop2}}{=} \pm J(v,u).
	\end{equation*}
Hence, when the bilinear form $J$ is symmetric or skew-symmetric, it is compatible with~\eqref{SubalgCond}.
\end{proof}

One may define color extensions of orthogonal, symplectic, and orthosymplectic Lie superalgebras, and so on.
A study in this direction, as well as the quantization of color Lie algebras and their affine analogue, is carried out in~\cite{zhang2026}.

Finally, we define a graded representation.

\begin{Definition}
	A graded representation $(\rho,V) $ of a color Lie algebra $\g$ in a $\Gamma$-graded vector space $V$ is a color Lie algebra homomorphism $ \rho\colon \g \to \gl(V,\omega)$ of degree zero:
	\begin{equation*}
		\rho(\llbracket X, Y \rrbracket) = \llbracket \rho(X), \rho(Y) \rrbracket,
		\qquad X, Y \in \g.
	\end{equation*}
\end{Definition}


\section{Graded Casimir elements and graded central extensions} \label{SEC:Method}

\subsection[Invariant bilinear form on color Lie algebras and 2nd order Casimir elements]{Invariant bilinear form on color Lie algebras\\ and 2nd order Casimir elements}

Let $\g$ be a color Lie algebra defined for a commutative factor $\omega$ on an Abelian group $\Gamma$.
We begin with some definitions.

\begin{Definition}\label{DEF:gcenter}
	The graded center $Z(\g)$ of $\g$ is the set of elements which graded-commute with all elements of $\g$:
	\begin{equation*}
		Z(\g) = \{ X \in \g \mid \llbracket X, Y \rrbracket = 0, \ \forall Y \in \g \}.
	\end{equation*}
\end{Definition}

Let $T(\g) := \sum_{r=0}^{\infty	} \g^{\otimes r} $ be the tensor algebra over~$\g$.
$T(\g)$ has $ \mathbb{Z}_+ \times \Gamma$-grading. The $\mathbb{Z}_+$ grading is obvious and the $\Gamma$-grading is given by
\begin{equation*}
	T(\g) = \sum_{\alpha\in \Gamma} T(\g)_{\alpha},
	\qquad
	T(\g)_{\alpha} = \sum_{r=0}^{\infty}\biggl( \sum_{\alpha_1+ \cdots + \alpha_r = \alpha} \g_{\alpha_1} \otimes \g_{\alpha_2} \otimes \cdots \otimes \g_{\alpha_r} \biggr).
\end{equation*}
Let $I(\g)$ be the two-sided ideal of $T(\g)$ generated by
\begin{equation*}
	X \otimes Y - \omega(\alpha,\beta) Y \otimes X - \llbracket X, Y \rrbracket, \qquad X \in \g_{\alpha},\quad Y \in \g_{\beta}.
\end{equation*}
The universal enveloping algebra of $\g$ is defined by $U(\g) = T(\g)/I(\g).$
Obviously, $U(\g)$ is a unital associative $\Gamma$-graded algebra.

\begin{Definition}
	Graded Casimir elements of $\g$ are elements of the graded center of $U(\g)$.
	Namely, if $C \in U(\g)$ is a Casimir element of degree $\gamma$, $ C X = \omega(\gamma, \alpha) XC $ for all homogeneous elements $ X \in U(\g) $ of degree~$\alpha$.
\end{Definition}

We denote an ordered basis of $\g_{\alpha}$ by $ X_{\alpha_i}$, $1 \leq i \leq \dim \g_{\alpha}$, and write the defining relations of~$\g$~as
\begin{align}
	& 	\llbracket X_{\alpha_i}, X_{\beta_j} \rrbracket = \sum_{\gamma_k} f_{\alpha_i, \beta_j}^{\quad\ \gamma_k} X_{\gamma_k}, \qquad f_{\alpha_i, \beta_j}^{\quad\ \gamma_k} \in \mathbb{C}, \label{sumgammak}
	\\
	& f_{\alpha_i, \beta_j}^{\quad\ \gamma_k} = 0 \qquad \text{if} \quad \alpha + \beta \neq \gamma, \nn
\end{align}
where $\sum_{\gamma_k}$ is a shorthand notation for $\sum_{\gamma \in \Gamma} \sum_{k=1}^{\dim \g_{\gamma}}$.
The relation
\[
\llbracket X_{\alpha_i}, X_{\beta_j} \rrbracket = -\omega(\alpha,\beta) \llbracket X_{\beta_j}, X_{\alpha_i} \rrbracket
\]
 is equivalent to the relation
\begin{equation}
	f_{\alpha_i, \beta_j}^{\quad\ \gamma_k} = -\omega(\alpha,\beta) f_{\beta_j,\alpha_i}^{\quad\ \gamma_k}.
	\label{fskew}
\end{equation}
Let $(\rho,V)$ be an irreducible graded representation of $\g$ in a $\Gamma$-graded vector space $V$.
If a~degree~$\mu$ element $ M^{\mu} \in \operatorname{End}(V)$ graded-commutes with $\rho(X)$ for all $ X \in \g$, i.e.,
\begin{equation*}
	\llbracket M^{\mu}, \rho(X) \rrbracket = 0, \qquad \forall X \in \g,
\end{equation*}
then we call $M^{\mu}$ a \textit{commutant} of degree $\mu$.
Clearly, the degree zero commutant is (a multiple of) the identity.

\begin{Proposition} \label{PROP:BF}
Assume $M^{-\mu}$ is a commutant of degree $-\mu$.
	We define a bilinear form $\eta^{\mu}\colon \g \times \g \to \mathbb{C}$ by
	\begin{equation*}
		\eta^{\mu}(X_{\alpha_i}, X_{\beta_j}) = \operatorname{ctr}( \rho(X_{\alpha_i}) M^{-\mu} \rho(X_{\beta_j}) ), %\label{formdef}
	\end{equation*}
	which is also denoted by $\eta^{\mu}_{\alpha_i,\beta_j}$. This bilinear form satisfies the following properties:
	\begin{enumerate}\itemsep=0pt
		\item[$(1)$] $ \eta^{\mu}_{\alpha_i, \beta_j} = 0$ if $ \alpha + \beta \neq \mu$,
		%
		\item[$(2)$] $ \eta^{\mu}_{\beta_j,\alpha_i} = \omega(\mu,\mu) \omega(\alpha,\beta) \eta^{\mu}_{\alpha_i, \beta_j} $,
		%
		\item[$(3)$] $ \eta^{\mu}(\llbracket X_{\alpha_i}, X_{\beta_j} \rrbracket, X_{\gamma_k} ) = \omega(\mu,\beta) \eta^{\mu}(X_{\alpha_i}, \llbracket X_{\beta_j}, X_{\gamma_k} \rrbracket) $.
	\end{enumerate}
	Relation $(3)$ implies the $\g$-invariance of $\eta^{\mu}$.
\end{Proposition}

\begin{proof}
	To simplify the notations, we set $ A := \rho(X_{\alpha_i})$, $B := \rho(X_{\beta_j})$, $C := \rho(X_{\gamma_k})$.

(1) By the definition of the color trace, $\operatorname{ctr} X = 0 $ if $\deg X \neq 0.$
		The statement immediately follows from $ \deg (AM^{-\mu}B) = \alpha + \beta - \mu$.

(2) The quantity $\eta^{\mu}_{\beta_j,\alpha_i} = \operatorname{ctr}(B M^{-\mu} A) $ is computed, using \eqref{ctrperm} and $ \llbracket M^{-\mu}, A \rrbracket = 0, $ as follows:
		\begin{align*}
			\eta^{\mu}_{\beta_j,\alpha_i} &
			\stackrel{\eqref{ctrperm}}{=} \omega(\beta,-\mu+\alpha)\, \operatorname{ctr}(M^{-\mu} A B)
			= \omega(\beta,-\mu+\alpha) \omega(-\mu,\alpha)\, \operatorname{ctr}(A M^{-\mu} B)
			\nn \\
			&\, \, = \omega(\beta,-\mu+\alpha) \omega(-\mu,\alpha)\, \eta^{\mu}_{\alpha_i, \beta_j}.
		\end{align*}
		Recalling that $ \eta^{\mu}_{\alpha_i, \beta_j} \neq 0 $ only if $ \alpha+\beta = \mu,$
		it is straightforward to verify the identity
\[
\omega(\beta,-\mu+\alpha) \omega(-\mu,\alpha) = \omega(\mu,\mu) \omega(\alpha,\beta).
\]

(3) By definition, we can write
		\begin{align*}
			\eta^{\mu}(\llbracket X_{\alpha_i}, X_{\beta_j} \rrbracket, X_{\gamma_k} ) =
			\operatorname{ctr}(\llbracket A, B \rrbracket M^{-\mu}C)
			=\operatorname{ctr}\big( (AB - \omega(\alpha,\beta) BA) M^{-\mu}C \big).
		\end{align*}
		Using \eqref{ctrperm} and \eqref{omegaprop2}, we obtain
		\begin{align*}
			\eta^{\mu}(\llbracket X_{\alpha_i}, X_{\beta_j} \rrbracket, X_{\gamma_k} )
			&= \operatorname{ctr}\big( \omega(\beta,-\mu) A M^{-\mu} BC - \omega(\alpha,\beta) \omega(\beta,\alpha-\mu+\gamma) A M^{-\mu}CB \big)
			\nn \\
			&= \omega(\mu,\beta) \operatorname{ctr}\big( A M^{-\mu} (BC - \omega(\beta,\gamma) CB)) \big)
			\nn \\
			&= \omega(\mu,\beta) \operatorname{ctr} (A M^{-\mu} \llbracket B, C \rrbracket) = \omega(\mu,\beta) \eta^{\mu}(X_{\alpha_i}, \llbracket X_{\beta_j}, X_{\gamma_k} \rrbracket).
		\end{align*}
This completes the proof.
\end{proof}

In terms of the structure constants $f_{\alpha_i, \beta_j}^{\quad\ \gamma_k}$, the $\g$-invariance of $\eta^{\mu}$ takes the following form:
\begin{equation}
	\sum_{\nu_{\ell}}
	f_{\alpha_i \beta_j}^{\quad \ \nu_{\ell}} \eta^{\mu}_{\nu_{\ell}\gamma_k}
	=
	\omega(\mu,\beta)
	\sum_{\nu_{\ell}}
	f_{\beta_j \gamma_k}^{\quad \ \nu_{\ell}} \eta^{\mu}_{\alpha_i \nu_{\ell}},
	\label{Invf}
\end{equation}
where we continue to use the shorthand notation introduced in~\eqref{sumgammak}.
If $\eta^{\mu}$ is nondegenerate, we introduce its inverse and denote it by $ \eta_{\mu}$:
\begin{align*}
	 &\eta_{\mu}^{\alpha_i \beta_j}=\eta_{\mu}(X_{\alpha_i}, X_{\beta_j}) = 0,\qquad \alpha + \beta \neq \mu,
	\qquad
	 \sum_{\gamma_k } \eta_{\mu}^{\alpha_i \gamma_k} \eta^{\mu}_{\gamma_k \beta_j}=\sum_{\gamma_k}
	\eta^{\mu}_{\beta_j \gamma_k} \eta_{\mu}^{\gamma_k \alpha_i} = \delta^{\alpha_i}_{\beta_j}.
\end{align*}
An equivalent expression to \eqref{Invf} in terms of $\eta_{\mu}$ is given by
\begin{equation}
	\sum_{\nu_{\ell}}
	\eta_{\mu}^{\alpha_i \nu_{\ell}} f_{\nu_{\ell} \beta_j}^{\quad \ \gamma_k}
	= \omega(\mu,\beta)
	\sum_{\nu_{\ell}}
	\eta_{\mu}^{\nu_{\ell} \gamma_k} f_{\beta_j \nu_{\ell}}^{\quad\ \alpha_i}.
	\label{Invf2}
\end{equation}

\begin{Theorem}
	A second order Casimir element of degree $\mu$ is given by
	\begin{equation*}
		C_{\mu} = \sum_{\alpha_i,\beta_j} \eta_{\mu}^{\alpha_i \beta_j} X_{\alpha_i} X_{\beta_j}.
	\end{equation*}
\end{Theorem}

\begin{proof}
	It is sufficient to prove $ \llbracket X_{\alpha_i}, C_{\mu} \rrbracket = 0. $
	It is proved by straightforward computation:
	\begin{align*}
		\llbracket X_{\alpha_i}, C_{\mu} \rrbracket &= \llbracket X_{\alpha_i}, \sum_{\beta_j,\gamma_k} \eta_{\mu}^{\beta_j \gamma_k} X_{\beta_j} X_{\gamma_k} \rrbracket
		\nn\\
		&= \sum_{\beta_j,\gamma_k,\nu_{\ell}} \eta_{\mu}^{\beta_j \gamma_k}
		\big( f_{\alpha_i \beta_j}^{\quad \ \nu_{\ell}} X_{\nu_{\ell}} X_{\gamma_k} + \omega(\alpha,\beta) f_{\alpha_i \gamma_k}^{\quad\ \nu_{\ell}} X_{\beta_j} X_{\nu_{\ell}} \big)
		\nn\\
		&=
		\sum_{\beta_j,\gamma_k,\nu_{\ell}}
		\big( \eta_{\mu}^{\beta_j \gamma_k} f_{\alpha_i \beta_j}^{\quad \ \nu_{\ell}} + \omega(\alpha,\nu) \eta_{\mu}^{\nu_{\ell} \beta_j} f_{\alpha_i \beta_j}^{\quad \ \gamma_k} \big)
		X_{\nu_{\ell}} X_{\gamma_k}
		\nn \\
		&\stackrel{\eqref{fskew}}{=}
		\sum_{\beta_j,\gamma_k,\nu_{\ell}}
		\big( \eta_{\mu}^{\beta_j \gamma_k} f_{\alpha_i \beta_j}^{\quad \ \nu_{\ell}} - \omega(\alpha,\nu) \omega(\alpha,\beta)\, \eta_{\mu}^{\nu_{\ell} \beta_j} f_{\beta_j \alpha_i}^{\quad \ \gamma_k} \big)
		X_{\nu_{\ell}} X_{\gamma_k}
		\nn\\
		&\stackrel{\eqref{Invf2}}{=}
		\sum_{\beta_j,\gamma_k,\nu_{\ell}} \eta_{\mu}^{\nu_{\ell} \beta_j} f_{\beta_j \alpha_i}^{\quad \ \gamma_k} \big( \omega(\alpha,\mu) - \omega(\alpha,\nu+\beta) \big)
		X_{\nu_{\ell}} X_{\gamma_k}
		 \stackrel{\nu+\beta=\mu}{=} 0.\tag*{\qed}
	\end{align*}\renewcommand{\qed}{}
\end{proof}

\section[Z\_2\^{}2-graded extension of osp(m|2n)]{$\boldsymbol{\Z2}$-graded extension of $\boldsymbol{\osp(m|2n)}$} \label{SEC:osp}

In this section, we consider a particular $\Z2$-graded extension of the Lie superalgebra $\osp(m|2n)$ and show that the color Lie algebra has a commutatnt of degree $11.$
The grading group $\Gamma$ and the commutative factor $\omega(\alpha,\beta) = (-1)^{\alpha_1 \beta_1 + \alpha_2\beta_2}$ are the same as in Section~\ref{SEC:qn}.

\subsection{Definition, commutants and basis}

We define the $\Z2$-graded color algebra, denoted by \Zosp, as a subalgebra of algebra $ \gl(m,m|2n,2n)$.
An element of $\mathfrak{gl}(m,m|2n,2n)$ is given by a $(2m+4n) \times (2m+4n)$ matrix with a $4 \times 4$ block structure, where each block has a definite $\Z2$-degree. We consider a particular class of elements of $\mathfrak{gl}(m,m|2n,2n)$ of the form
\begin{equation}
	X =
	\begin{pNiceMatrix}[first-row][first-col]
		& m & m & 2n & 2n \\
		m & A^{00} & A^{11} & A^{01} & A^{10}
		\\
		m & A^{11} & A^{00} & A^{10} & A^{01}
		\\
		2n & C^{01} & C^{10} & C^{00} & C^{11}
		\\
		2n & C^{10} & C^{01} & C^{11} & C^{00}
	\end{pNiceMatrix},
	\label{doubledGL}
\end{equation}
where the superscripts indicate the $\Z2$-degree of each block, and each block matrix appears twice.

The bilinear form defining \Zosp\ is given by a $ 4 \times 4$ block matrix of the same size as~$X$:
\begin{equation*}
	J =
	\begin{pmatrix}
		{\rm i}B_m & 0 & 0 & 0
		\\
		0 & {\rm i}B_m & 0 & 0
		\\
		0 & 0 & G_{2n} & 0
		\\
		0 & 0 & 0 & G_{2n}
	\end{pmatrix}.
\end{equation*}
Here,
\begin{align*}
	&B_m=
	\begin{pmatrix}
		0 & \mathbb{I}_{\ell} \\ \mathbb{I}_{\ell} & 0
	\end{pmatrix}, \quad m = 2\ell,
	\qquad
	B_m =
	\begin{pmatrix}
		0 & \mathbb{I}_{\ell} & 0 \\ \mathbb{I}_{\ell} & 0 & 0
		\\
		0 & 0 & 1
	\end{pmatrix}, \quad m = 2\ell+1,
	\notag \\
	&G_{2n}=
	\begin{pmatrix}
		0 & -\mathbb{I}_{n} \\ \mathbb{I}_{n} & 0
	\end{pmatrix},
\end{align*}
where $ \mathbb{I}_{\ell} $ denotes the $\ell \times \ell$ identity matrix.

To define \Zosp, we impose the following conditions on the matrix $X$ given in \eqref{doubledGL}:
\begin{alignat}{3}
&	X^{00}J + J \bigl(X^{00}\bigr)^{\mathsf{T}}= 0,\qquad &&
	X^{11}J + J \bigl(X^{11}\bigr)^{\mathsf{T}}= 0,&
	\notag\\
&	X^{01}J + {\rm i}J \bigl(X^{01}\bigr)^{\mathsf{T}}= 0,
	\qquad &&
	X^{10}J -{\rm i} J \bigl(X^{10}\bigr)^{\mathsf{T}}= 0.& \label{ospCondition}
\end{alignat}
Here $X^{\alpha}$, with $\alpha \in \mathbb{Z}_2$, denotes the matrix in~\eqref{doubledGL} whose entries are nonzero only in the blocks of $\mathbb{Z}_2^2$-degree $\alpha$; e.g.,
\begin{equation*}
	X^{11} =
	\begin{pmatrix}
		0 & A^{11} & 0 & 0
		\\
		A^{11} & 0 & 0 & 0
		\\
		0 & 0 & 0 & C^{11}
		\\
		0 & 0 & C^{11} & 0
	\end{pmatrix}
\end{equation*}
and $ X^{\mathsf{T}} $ denotes the transpose of $X.$
It is straightforward to check by direct computation that if $X^{\alpha}$ and $Y^{\beta}$ satisfy the conditions in \eqref{ospCondition}, then so does $\llbracket X^{\alpha}, Y^{\beta} \rrbracket$.

The conditions in \eqref{ospCondition} yield the following constraints on the matrix blocks:
\begin{alignat}{3}
	&A^{00} B_m+ B_m \bigl(A^{00}\bigr)^{\mathsf{T}} = 0, \qquad && A^{11} B_m+ B_m \bigl(A^{11}\bigr)^{\mathsf{T}} = 0,&
	\notag\\
&	C^{00}G_{2n}+ G_{2n}\bigl(C^{00}\bigr)^{\mathsf{T}} = 0,\qquad&& C^{11} G_{2n}+ G_{2n}\bigl(C^{11}\bigr)^{\mathsf{T}} = 0,&
	\notag \\
&	C^{01}= -G_{2n}\bigl(A^{01}\bigr)^{\mathsf{T}} B_m, \qquad && C^{10}= G_{2n}\bigl(A^{10}\bigr)^{\mathsf{T}} B_m.&
	\label{StrOfX}
\end{alignat}
It follows that
\begin{equation*}
	A^{00}, A^{11} \in \so(m), \qquad C^{00}, C^{11} \in \sp(2n)
\end{equation*}
and $ A^{01}$, $A^{10}$ are arbitrary $ m \times 2n$ matrices.
Therefore, we see that
\[
\dim \Z2\text{-}\osp(m|2n) = (m+2n)^2-m+2n,
\]
which is twice the dimension of the Lie superalgebra $\osp(m|2n)$.

We comment on preceding works on $\Z2$-graded extensions of $\osp(m|2n)$ in turn:
\begin{itemize}\itemsep=0pt
	\item[(1)] Rittenberg and Wyler introduced a $\Z2$-graded extension of the Lie supergroup $Osp(m|2n)$ and define a $\Z2$-graded $\osp(m|2n)$ as its infinitesimal transformation \cite{rw2}.
	Algebraic definitions of a $\Z2$-graded extension of $\osp(m|2n)$ can be found in \cite{GrJav, JaYangWyb,SVdJ2018,SVdJ2024}.
	However, none of these works discuss the particular subclass defined by \eqref{doubledGL}.
	\item[(2)] A $\Z2$-graded extension of the type defined in \eqref{doubledGL}, which corresponds to $\Z2$-$\osp(1|2)$ in the present notation, first appeared as a $\Z2$-graded extension of conformal quantum mechanics~\cite{Aizawa_2020}.
	The same algebra was also obtained via a module construction of color Lie algebras~\cite{LU20231}.
	It was pointed out in \cite{AiSe} that $\Z2$-$\osp(1|2)$ has a $11$-graded commutant.
	\item[(3)] The lowest weight representations of $\Z2$-$\osp(1|2)$ are classified in \cite{AmaAi} (see \cite{AlHoSe} for other representations of $\Z2$-$\osp(1|2)$).
\end{itemize}

Now we return to the search for commutants of \Zosp.

\begin{Proposition} \label{PROP:comm}
	There exists a unique $($up to constant multiple$)$ $11$-graded commutant of \Zosp, which is given by
	\begin{equation}
		M^{11} =
		\begin{pNiceArray}{cc|cc}
			0 & \mathbb{I}_m & &
			\\
			\mathbb{I}_m & 0 & &
			\\ \hline
			& & 0 & -\mathbb{I}_{2n}
			\\
			& & -\mathbb{I}_{2n} & 0
		\end{pNiceArray}.
		\label{M11osp}
	\end{equation}
	On the other hand, there exist no $01$- or $10$-graded commutants.
\end{Proposition}

\begin{proof}
	$M^{11}$ must be of the form
	\begin{equation*}
		M^{11} =
		\begin{pNiceArray}{cc|cc}
			0 & N_1 & &
			\\
			N_2 & 0 & &
			\\ \hline
			& & 0 & N_3
			\\
			& & N_4 & 0
		\end{pNiceArray}.
		%\label{M11}
	\end{equation*}
	The condition $[X^{00}, M^{11}] = 0$ yields the following constraints on $N_k$:
	\begin{equation*}
		\bigl[A^{00},N_1\bigr] = \bigl[A^{00},N_2\bigr] = 0, \qquad \bigl[C^{00},N_3\bigr] = \bigl[C^{00},N_4\bigr] = 0.
	\end{equation*}
	Recalling that $ A^{00} \in \so(m) $ and $ C^{00} \in \sp(2n)$, we obtain, by Schur's lemma,
	\begin{equation*}
		N_1 = \alpha_1 \mathbb{I}_m, \qquad N_2 = \alpha_2 \mathbb{I}_m, \qquad
		N_3 = \alpha_3 \mathbb{I}_{2n}, \qquad N_4 = \alpha_4 \mathbb{I}_{2n}, \qquad \alpha_k \in \mathbb{C}.
	\end{equation*}
	The condition $[X^{11}, M^{11}] = 0$ further implies that $\alpha_1 = \alpha_2$ and $\alpha_3 = \alpha_4$. Finally, the condition $\bigl\{ X^{01}, M^{11}\bigr\} = 0$ yields $ \alpha_3 = -\alpha_1$, and this relation is consistent with $\bigl\{ X^{10}, M^{11}\bigr\} = 0$.
	Thus, $M^{11}$ is given by~\eqref{M11osp}.

	We can search for $M^{01}$ in a similar way. We write
\begin{equation*}
	M^{01} =
	\begin{pNiceArray}{cc|cc}
		& & N_1 & 0
		\\
		& & 0 & N_2 \\ \hline
		N_3 & 0 & &
		\\
		0 & N_4 & &
	\end{pNiceArray}
\end{equation*}
and the condition $[X^{00}, M^{01}] = 0$ yields the constraints
\[
	A^{00} N_1 = N_1 C^{00}, \qquad A^{00} N_2 = N_2 C^{00}, \qquad
	C^{00} N_3 = N_3 A^{00}, \qquad C^{00} N_4 = N_4 A^{00}.
\]
If $ A^{00} = 0 $ and $ C^{00} \neq 0$, then $ N_k \equiv 0$, and hence $ M^{01} = 0.$
In the same way, one finds that $M^{10} = 0$.
\end{proof}

From Proposition~\ref{PROP:comm}, we see that \Zosp\ has $00$- and $11$-graded 2nd order Casimir elements.
To derive explicit formulas for them, we first need to specify a basis of \Zosp.
This is also helpful to understand structures of the color Lie algebra.
Recall that $ X \in $ \Zosp\ has a $ 4 \times 4 $ block structure, see \eqref{doubledGL}, and that each block further has the following substructure.
In the case of $ m = 2\ell+1$,
\begin{equation*}
	A^{00} =
	\begin{pmatrix}
		L & Y & x
		\\
		W & -L^{\mathsf{T}} & y
		\\
		-y^{\mathsf{T}} & -x^{\mathsf{T}} & 0
	\end{pmatrix}
	\in \so(2\ell+1), \qquad
	C^{00} =
	\begin{pmatrix}
		\tilde{L} & \tilde{Y} \\ \tilde{W} & -\tilde{L}^{\mathsf{T}}
	\end{pmatrix}
	\in \sp(2n),
\end{equation*}
where $ L \in \gl(\ell)$ and $ Y$, $W $ are skew-symmetric $\ell \times \ell $ matrices, and $ x, y \in \mathbb{C}^{\ell}$,
while $\tilde{L} \in \gl(n)$ and $ \tilde{Y}$, $\tilde{W}$ are symmetric $ n \times n $ matrices.
The block structures of $A^{01}$ and $A^{10}$ are induced from this structure:
\begin{equation*}
	A^{01} =
	\begin{pmatrix}
		P & Q \\ R & S \\
		a & b
	\end{pmatrix},
	\qquad
	A^{10} =
	\begin{pmatrix}
		\tilde{P} & \tilde{Q} \\ \tilde{R} & \tilde{S} \\
		\tilde{a} & \tilde{b}
	\end{pmatrix},
\end{equation*}
where $ P$, $Q$, $\tilde{P}$ and $\tilde{Q} $ are $ \ell \times n $ matrices, and $ a$, $b$, $\tilde{a} $ and $\tilde{b} $ are $n$-dimensional row vectors.
For $m = 2\ell$, the corresponding changes are obvious.
Then, \eqref{StrOfX} determines $ C^{01}$ and $ C^{10}$ as follows:%
\begin{equation*}
	C^{01}
	=
	\begin{pmatrix}
		S^{\mathsf{T}} & Q^{\mathsf{T}} & b^{\mathsf{T}} \\
		-R^{\mathsf{T}} & -P^{\mathsf{T}} & -a^{\mathsf{T}}
	\end{pmatrix},
	\qquad
	C^{10} =
	\begin{pmatrix}
		-\tilde{S}^{\mathsf{T}} & - \tilde{Q}^{\mathsf{T}} & -\tilde{b}^{\mathsf{T}}
		\\
		\tilde{R}^{\mathsf{T}} & \tilde{P}^{\mathsf{T}} & \tilde{a}^{\mathsf{T}}
	\end{pmatrix}.
\end{equation*}
Therefore, we introduce a concise notation for matrix units, by which the position of a nonvanishing entry can be easily read off.

Let $ e_{ab}$ be the matrix unit of size $(2m+4n)\times (2m+4n)$.
From now on, Latin indices take values in $\{ 1, 2, \dots, \ell= \mathrm{rank}\ \so(m) \}$, while Greek indices take values in $\{ 1, 2, \dots, n = \mathrm{rank}\ \sp(2n) \}$.
We~introduce the following notation for the matrix units whose nonvanishing entries lie in the first $\ell $ rows of the matrix $ X $ in \eqref{doubledGL}: For $ m = 2\ell$,
\begin{alignat*}{5}
	&E_{ij}:= e_{ij}, \qquad&& E_{ij'}:= e_{i\, \ell+j}, \qquad&& {E_i}^j:= e_{i\, m+j},\qquad&& {E_{i}}^{j'}:= e_{i\, m+\ell+j},&
	\notag\\
&	E_{i\nu}:=e_{i\, 2m+\nu}, \qquad && E_{i\nu'}:= e_{i\, 2m+n+\nu}, \qquad && {E_i}^{\nu}:= e_{i\, 2m+2n+\nu}, \qquad && {E_i}^{\nu'}:= e_{i\, 2m+3n+\nu},&
\end{alignat*}
and we need two more for $m = 2\ell+1$:
\begin{equation*}
	E_{im} := e_{i\, 2\ell+1}, \qquad {E_i}^m := e_{i,4\ell+2}.
\end{equation*}
The meaning of this notation is clear; $E_{ij}$ has its nonvanishing entry in the $\ell \times \ell$ block at the top left corner of $X$.
As the second index changes, the position of the nonvanishing entry shifts to the right.
The same rule can be applied to the first index (with shifts downward), allowing us to introduce a notation for the matrix units whose nonvanishing entries lie in the blocks: % indicated in \eqref{Etable}:
\begin{equation*}\renewcommand{\arraystretch}{1.3}
		\begin{pNiceArray}{cc:c|cc:c|cc|cc}[first-row,first-col]
			& \ell & \ell & 1 & \ell & \ell & 1 & n & n & n & n
			\\
			\ell & E_{ij} & E_{ij'} & E_{im} & {E_i}^{j} & {E_i}^{j'} & {E_{i}}^m & E_{i\nu} & E_{i\nu'} & {E_i}^{\nu} & {E_i}^{\nu'} \\
			\ell & E_{i'j} & E_{i'j'} & E_{i'm} & {E_{i'}}^{j} & {E_{i'}}^{j'} & {E_{i'}}^m & E_{i'\nu} & E_{i'\nu'} & {E_{i'}}^{\nu} & {E_{i'}}^{\nu'} \\ \hdashline
			%
			1 & E_{mj} & E_{mj'} & E_{mm} & {E_m}^{j} & {E_m}^{j'} & {E_{m}}^m & E_{m\nu} & E_{m\nu'} & {E_m}^{\nu} & {E_m}^{\nu'} \\ \hline
			%
			\ell &{E^i}_{j} & {E^i}_{j'} & {E^{i}}_{m} & E^{ij} & E^{ij'} & E^{im} & {E^i}_{\nu} & {E^i}_{\nu'} & E^{i\nu} & E^{i\nu'} \\
			%
			\ell &{E^{i'}}_{j} & {E^{i'}}_{j'} & {E^{i'}}_{m} & E^{i'j} & E^{i'j'} & E^{i'm} & {E^{i'}}_{\nu} & {E^{i'}}_{\nu'} & E^{i'\nu} & E^{i'\nu'} \\ \hdashline
			%
			1 & {E^m}_{j} & {E^m}_{j'} & {E^m}_{m} & E^{mj} & E^{mj'} & E^{mm} & {E^{m}}_{\nu} & {E^{m}}_{\nu'} & E^{m\nu} & E^{m\nu'} \\ \hline
			%
			n & E_{\mu j} & E_{\mu j'} & E_{\mu m} & {E_\mu}^{j} & {E_\mu}^{j'} & {E_{\mu}}^m & E_{\mu\nu} & E_{\mu\nu'} & {E_\mu}^{\nu} & {E_\mu}^{\nu'} \\
			%
			n & E_{\mu'j} & E_{\mu'j'} & E_{\mu' m} & {E_{\mu'}}^{j} & {E_{\mu'}}^{j'} & {E_{\mu'}}^m & E_{\mu'\nu} & E_{\mu'\nu'} & {E_{\mu'}}^{\nu} & {E_{\mu'}}^{\nu'} \\ \hline
			%
			n & {E^\mu}_{j} & {E^\mu}_{j'} & {E^{\mu}}_{m} & E^{\mu j} & E^{\mu j'} & E^{\mu m} & {E^\mu}_{\nu} & {E^\mu}_{\nu'} & E^{\mu\nu} & E^{\mu\nu'} \\
			%
			n & {E^{\mu'}}_{j} & {E^{\mu'}}_{j'} & {E^{\mu'}}_{m} & E^{\mu'j} & E^{\mu'j'} & E^{\mu' m} & {E^{\mu'}}_{\nu} & {E^{\mu'}}_{\nu'} & E^{\mu'\nu} & E^{\mu'\nu'}
		\end{pNiceArray}.
	%	\label{Etable}
	\end{equation*}
Note that the prime on $m$ is omitted for simplicity.
One advantage of this notation is the simplicity of the product: if the adjacent indices are not of the same type, the product vanishes. For example,
\begin{equation*}
	E_{ij'} E_{pq} = {E_i}^j E^{\mu \nu} = {E_i}^j {E_p}^q = 0, \qquad
	E_{ij}{E_p}^q = \delta_{jp} {E_i}^q.
\end{equation*}

Using this notation of the block structures described above, we choose the following basis of \Zosp\ in the case of $ m= 2\ell + 1$:
\begin{enumerate}\itemsep=0pt
\item[(i)] basis of degree $00$
\begin{align*}
	&T_{ij}=E_{ij}-E_{j'i'}+E^{ij}-E^{j'i'}, \qquad 1 \leq i, j \leq \ell, \text{ or } j=m,
	\notag\\
	&T_{ij'}=E_{ij'}-E_{ji'}+E^{ij'}-E^{ji'},\qquad 1 \leq i< j \leq \ell,
	\notag\\
	&T_{i'j}=E_{i'j}-E_{j'i}+E^{i'j}-E^{j'i},\qquad 1 \leq i< j \leq \ell, \text{ or } j=m,
	\notag\\
	&T_{\mu\nu}=E_{\mu\nu}-E_{\nu'\mu'}+E^{\mu\nu}-E^{\nu'\mu'}, \qquad 1 \leq \mu, \nu \leq n,
	\notag\\
	&T_{\mu\nu'}=
	\begin{cases}
		E_{\mu\nu'}+E_{\nu\mu'}+E^{\mu\nu'}+E^{\nu\mu'},& 1 \leq \mu<\nu \leq n,\\
		E_{\mu\mu'}+E^{\mu\mu'},&\mu=\nu,
	\end{cases}
	\notag\\
	&T_{\mu'\nu}=
	\begin{cases}
		E_{\mu'\nu}+E_{\nu'\mu}+E^{\mu'\nu}+E^{\nu'\mu},& 1 \leq \mu<\nu \leq n,\\
		E_{\mu'\mu}+E^{\mu'\mu}, & \mu=\nu.
	\end{cases}
\end{align*}
\item[(ii)] basis of degree $11$
\begin{align*}
	&{U_{i}}^{j}={E_{i}}^{j}-{E_{j'}}^{i'}+{E^{i}}_{j}-{E^{j'}}_{i'}, \qquad 1 \leq i, j \leq \ell, \text{ or } j=m,
	\notag\\
	&{U_{i}}^{j'}={E_{i}}^{j'}-{E_{j}}^{i'}+{E^{i}}_{j'}-{E^{j}}_{i'},\qquad 1 \leq i< j \leq \ell,
	\notag\\
	&{U_{i'}}^{j}={E_{i'}}^{j}-{E_{j'}}^{i}+{E^{i'}}_{j}-{E^{j'}}_{i},\qquad 1 \leq i< j \leq \ell, \text{ or } j=m,
	\notag\\
	&{U_{\mu}}^{\nu}={E_{\mu}}^{\nu}-{E_{\nu'}}^{\mu'}+{E^{\mu}}_{\nu}-{E^{\nu'}}_{\mu'},
	\notag\\
	&{U_{\mu}}^{\nu'}=
	\begin{cases}
		{E_{\mu}}^{\nu'}+{E_{\nu}}^{\mu'}+{E^{\mu}}_{\nu'}+{E^{\nu}}_{\mu'},& 1 \leq \mu<\nu \leq n,\\
		{E_{\mu}}^{\mu'}+{E^{\mu}}_{\mu'},& \mu=\nu,
	\end{cases}
	\notag\\
	&{U_{\mu'}}^{\nu}=
	\begin{cases}
		{E_{\mu'}}^{\nu}+{E_{\nu'}}^{\mu}+{E^{\mu'}}_{\nu}+{E^{\nu'}}_{\mu},& 1 \leq \mu<\nu \leq n,\\
		{E_{\mu'}}^{\mu}+{E^{\mu'}}_{\mu},& \mu=\nu.
	\end{cases}
\end{align*}
\item[(iii)] basis of degree $01$, ($1 \leq i \leq \ell$, $1 \leq \mu \leq n$)
\begin{align*}
	&\Lambda_{i\mu}=E_{i\mu}-E_{\mu'i'}+E^{i\mu}-E^{\mu'i'},
	\notag\\
	&\Lambda_{i\mu'}=E_{i\mu'}+E_{\mu i'}+E^{i\mu'}+E^{\mu i'},
	\notag\\
	&\Lambda_{i'\mu}=E_{i'\mu}-E_{\mu'i}+E^{i'\mu}-E^{\mu'i},
	\notag\\
	&\Lambda_{i'\mu'}=E_{i'\mu'}+E_{\mu i}+E^{i'\mu'}+E^{\mu i}.
\end{align*}
In the first and the second lines, we allow $ i =m$.
\item[(iv)] basis of degree $10$, ($1 \leq i \leq \ell$, $1 \leq \mu \leq n$)
\begin{align*}
&	{\Gamma_{i}}^{\mu}={E_{i}}^{\mu}+{E_{\mu'}}^{i'}+{E^{i}}_{\mu}+{E^{\mu'}}_{i'},
	\notag\\
&	{\Gamma_{i}}^{\mu'}={E_{i}}^{\mu'}-{E_{\mu}}^{i'}+{E^{i}}_{\mu'}-{E^{\mu}}_{i'},
	\notag\\
&	{\Gamma_{i'}}^{\mu}={E_{i'}}^{\mu}+{E_{\mu'}}^{i}+{E^{i'}}_{\mu}+{E^{\mu'}}_{i},
	\notag\\
&	{\Gamma_{i'}}^{\mu'}={E_{i'}}^{\mu'}-{E_{\mu}}^{i}+{E^{i'}}_{\mu'}-{E^{\mu}}_{i}.
\end{align*}
In the first and the second lines, we allow $ i =m$.
The basis for the case $m = 2\ell$ is obtained by eliminating the elements with $i=m $ and $j=m$.
\end{enumerate}


\subsection[Structure of Z\_2\^{}2-osp(m|2n) and 2nd order Casimir elements]{Structure of $\boldsymbol{\Z2}$-$\boldsymbol{\osp(m|2n)}$ and 2nd order Casimir elements}

The Cartan subalgebra of \Zosp\ is defined as the maximal commuting subalgebra of degree $00$.
Obviously, it is spanned by
\begin{alignat}{3}
&	H_i:= T_{ii} = E_{ii} - E_{i'i'} + E^{ii} - E^{i'i'}, \qquad&&
	 1 \leq i \leq \ell ,&
	\nn \\
	&H_{\mu}:= T_{\mu\mu} = E_{\mu\mu} - E_{\mu'\mu'} + E^{\mu\mu} - E^{\mu'\mu'},\qquad
	& & 1 \leq \mu \leq n,&
\end{alignat}
and has dimension $\ell+n$.

Using the Cartan subalgebra, we obtain the root space decomposition of $\g := $\Zosp\ as usual. Let $ \mathfrak{h} $ be the Cartan subalgebra of $\g$ and
\begin{equation*}
	\g^{(\lambda)} = \{ X \in \g \mid [H, X] = \lambda(H) X, \ \forall H \in \mathfrak{h} \}.
\end{equation*}
Note that, since $\deg(\mathfrak{h}) = 00$, the color Lie bracket is identical to the commutator.
First, observe that there exist zero roots of degree $11$:{\samepage
\begin{equation*}
	[H, {U_i}^i] = [H, {U_{\mu}}^{\mu}] = 0, \qquad \forall H \in \mathfrak{h}.
\end{equation*}
Hence, the zero root space has dimension $\ell+n$.}

Denote by $ \epsilon_i$, $\delta_{\mu}$ the basis of the dual space $\mathfrak{h}^*$ such that
\begin{equation*}
	\epsilon_i(H_j) = \delta_{ij}, \qquad
	\delta_{\mu}(H_{\nu}) = \delta_{\mu\nu}.
\end{equation*}
Then, it is not difficult to see that the basis of $\g$ introduced in the previous subsection yields a~basis of the root spaces, e.g.,
\begin{align*}
	&[H_i, T_{jk}]= (\delta_{ij} - \delta_{ik} ) T_{jk} = (\epsilon_j - \epsilon_k)(H_i) T_{jk},
	\nn \\
	&[H_i, {U_j}^k]=(\delta_{ij} - \delta_{ik} ) {U_j}^k = (\epsilon_j - \epsilon_k)(H_i) {U_j}^k,
	\nn \\
	&[H_i, \Lambda_{j\nu}]= \delta_{ij} \Lambda_{j\nu} = \epsilon_j(H_i) \Lambda_{j\nu},
	\nn \\
	&[H_i, {\Gamma_j}^{\nu}]= \delta_{ij}{\Gamma_j}^{\nu} = \epsilon_j(H_i) {\Gamma_j}^{\nu}.
\end{align*}
Repeating similar computations, we see that, apart from the zero root, the root system of $\g$ coincides with that of the ordinary Lie superalgebra~$\osp(m|2n)$.
However, each nonzero root space is two-dimensional and consists of basis elements of different degrees.
The root system of~$\g$, including the zero root, is summarized in the table below. For $m = 2\ell$,
	\begin{equation*}\renewcommand{\arraystretch}{1.2}
		\begin{NiceArray}{c|c||c|c}
			\lambda & \g^{(\lambda)} & \lambda & \g^{(\lambda)}
			\\ \hline
			0 & U_i^{\ i}, \quad U_{\mu}^{\ \mu} & \epsilon_i - \epsilon_j & T_{ij}, \quad U_j^{\ j}
			\\
			\epsilon_i+\epsilon_j & T_{ij'} \quad U_i^{\ j'} & -\epsilon_i - \epsilon_j & T_{j'k}, \quad U_{j'}^{\ k}
			\\
			\delta_{\mu}-\delta_{\nu} & T_{\mu\nu}, \quad U_{\mu}^{\ \nu} & &
			\\
			\delta_{\mu}+\delta_{\nu} & T_{\mu\nu'}, \quad U_{\mu}^{\ \nu'} & -\delta_{\mu}-\delta_{\nu} & T_{\mu' \nu}, \quad U_{\mu'}^{\ \;\nu}
			\\
			2\delta_{\mu} & T_{\mu\mu'},\quad U_{\mu}^{\ \mu'} & -2\delta_{\mu} & T_{\mu'\mu},\quad U_{\mu'}^{\ \;\mu}
			\\ \hline
			\epsilon_i - \delta_{\mu} & \Lambda_{i\mu}, \quad \Gamma_i^{\ \mu} &
			-\epsilon_i+ \delta_{\mu} & \Lambda_{i'\mu'}, \quad \Gamma_{i'}^{\ \; \mu'}
			\\
			\epsilon_i + \delta_{\mu} & \Lambda_{i\mu'}, \quad \Gamma_i^{\mu'} &
			-\epsilon_i - \delta_{\mu} & \Lambda_{i'\mu}, \quad \Gamma_{i'}^{\ \;\mu}
		\end{NiceArray}
	\end{equation*}
	For $ m = 2\ell + 1$, we have additional roots
	\begin{equation*}\renewcommand{\arraystretch}{1.2}
		\begin{NiceArray}{c|c||c|c}
			\lambda & \g^{(\lambda)} & \lambda & \g^{(\lambda)}
			\\ \hline
			\epsilon_i & T_{im}, \quad U_{i}^{\ m} & -\epsilon_i & T_{i'm}, \quad U_{i'}^{\ m}
			\nn \\ \hline
			\delta_{\mu} & \Lambda_{m\mu'}, \quad \Gamma_{m}^{\ \;\mu'} & -\delta_{\mu} & \Lambda_{m\mu}, \quad \Gamma_{m}^{\ \;\mu}
		\end{NiceArray}
	\end{equation*}


Let us now turn to the normalized invariant forms on $\g$ defined by
\begin{equation*}
	\eta^{00}(X,Y) = \frac{1}{4}\operatorname{ctr}(XY),
	\qquad
	\eta^{11}(X,Y) = \frac{1}{4}\operatorname{ctr}\bigl(XM^{11}Y\bigr),
	\qquad X, Y \in \g.
\end{equation*}
The symmetry properties of these bilinear forms and the degrees $\deg(X)$, $\deg(Y)$ for which they take nonvanishing values are summarized as follows:
\begin{equation*}
	\begin{NiceArray}{c|cc|c}
		& \deg(X) & \deg(Y) & \text{symmetry}
		\\ \hline
		\Block{4-1}{\eta^{00}(X,Y)} & 00 & 00 & \Block{2-1}{\text{symmetric}}
		\\
		& 11 & 11 &
		\\ \Hline[start=2,end=4]
		& 01 & 01 & \Block{2-1}{\text{skew-symmetric}}
		\\
		& 10 & 10 &
		\\ \hline
		\Block{4-1}{\eta^{11}(X,Y)} & 00 & 11 & \Block{4-1}{\text{symmetric}}
		\\
		& 11 & 00 &
		\\
		& 01 & 10 &
		\\
		& 10 & 01 &
	\end{NiceArray}
\end{equation*}

Using the fact that the $11$-graded commutant \eqref{M11osp} can be written as
\[
	M^{11} = \sum_{i}\bigl({E_i}^i + {E_{i'}}^{i'} + {E^i}_i + {E^{i'}}_{i'}\bigr)
	-\sum_{\mu} \bigl({E_{\mu}}^{\mu} + {E_{\mu'}}^{\mu'} + {E^{\mu}}_{\mu} + {E^{\mu'}}_{\mu'}\bigr),
\]
it is easy to verify that $M^{11}$ maps the $11$- and $10$-graded basis of $\g$ to those of degree $00$ and $01$, respectively:
\begin{alignat*}{4}
&	M^{11} {U_i}^j= T_{ij}, \qquad&& M^{11} {U_i}^{j'} = T_{ij'},\qquad&& M^{11} {U_{i'}}^j= T_{i'j},&
	\notag \\
&	M^{11} {U_{\mu}}^{\nu}= -T_{\mu\nu}, \qquad && M^{11} {U_{\mu}}^{\nu'}= -T_{\mu\nu'}, \qquad && M^{11} {U_{\mu'}}^{\nu}= -T_{\mu'\nu},&
\end{alignat*}
and{\samepage
\begin{equation*}
	M^{11}{\Gamma_i}^{\mu} = \Lambda_{i\mu}, \qquad M^{11}{\Gamma_i}^{\mu'} = \Lambda_{i\mu'}, \qquad M^{11}{\Gamma_{i'}}^{\mu} = \Lambda_{i'\mu}, \qquad M^{11}{\Gamma_{i'}}^{\mu'} = \Lambda_{i'\mu'}
\end{equation*}
Therefore, the computation of $\eta^{11}$ reduces to that of $\eta^{00}$.}

Now the computation of $\eta^{00}$ and $\eta^{11}$ is straightforward. Here we list the nonvanishing components for $m = 2\ell + 1$. For $m = 2\ell$, the components corresponding to the basis elements with index $m$ do not exist. We note that additional nonvanishing components can be obtained from the symmetry properties of the bilinear forms; for simplicity, we do not list them.
Components corresponding to degree $00$ and $11$ basis:
\begin{align}
&	\eta^{00}(T_{ij},T_{pq})= \eta^{00}({U_i}^j, {U_p}^q) = \eta^{11}(T_{ij},{U_p}^q) = \delta_{iq} \delta_{jp},
	\notag\\
&	\eta^{00}(T_{ij'}, T_{p'q})= \eta^{00}\bigl({U_i}^{j'}, {U_{p'}}^{q}\bigr) = \eta^{11}(T_{ij'},{U_{p'}}^q) = \eta^{11}\bigl(T_{i'j},{U_p}^{q'}\bigr) = -\delta_{ip} \delta_{jq},
	\notag \\
	&\eta^{00}(T_{im}, T_{j'm})= \eta^{00}({U_i}^m,{U_{j'}}^m) = \eta^{11}(T_{im},{U_{j'}}^m) = \eta^{11}(T_{i'm}, {U_j}^m) = -\delta_{ij},
	\notag \\
	&\eta^{00}(T_{\mu\nu}, T_{\rho\sigma})= \eta^{00}({U_{\mu}}^{\nu},{U_{\rho}}^{\sigma}) = -\eta^{11}(T_{\mu\nu},{U_{\rho}}^{\sigma})= -\delta_{\mu\sigma} \delta_{\nu\rho},
	\notag\\
	&\eta^{00}(T_{\mu\nu'}, T_{\rho'\sigma})= \eta^{00}\bigl({U_{\mu}}^{\nu'},{U_{\rho'}}^{\sigma}\bigr) = -\eta^{11}(T_{\mu\nu'},{U_{\rho'}}^{\sigma}) =
	-\eta^{11}\bigl(T_{\mu'\nu},{U_{\rho}}^{\sigma'}\bigr)
	\notag\\
	&\hphantom{\eta^{00}(T_{\mu\nu'}, T_{\rho'\sigma})}{}
=
	\begin{cases}
		-\delta_{\mu\rho} \delta_{\nu\sigma}, & \mu < \nu, \rho < \sigma
		\\[5pt]
		-\frac{1}{2} \delta_{\mu\rho}, & \mu = \nu, \rho=\sigma,
	\end{cases}
	\label{BFosp1}
\end{align}
and $01$, $10$ basis:
\begin{align}
&	\eta^{00}(\Lambda_{i\mu}, \Lambda_{j'\nu'})= -\eta^{00}(\Lambda_{i\mu'}, \Lambda_{j'\nu}) = -\eta^{00}\bigl({\Gamma_i}^{\mu}, {\Gamma_{j'}}^{\nu'}\bigr) = \eta^{00}({\Gamma_i}^{\mu'}, {\Gamma_{j'}}^{\nu}) = \delta_{ij}\delta_{\mu\nu},
	\notag\\
&	\eta^{00}(\Lambda_{m\mu}, \Lambda_{m\nu'})= -\eta^{00}\bigl({\Gamma_m}^{\mu}, {\Gamma_m}^{\nu'}\bigr) = \delta_{\mu\nu},
	\notag\\
&	\eta^{11}\bigl(\Lambda_{i\mu},{\Gamma_{j'}}^{\nu'}\bigr)= -\eta^{11}(\Lambda_{i\mu'},{\Gamma_{j'}}^{\nu}) = \eta^{11}\bigl(\Lambda_{i'\mu},{\Gamma_{j}}^{\nu'}\bigr) = -\eta^{11}(\Lambda_{i'\mu'},{\Gamma_j}^{\nu}) = \delta_{ij} \delta_{\mu\nu},
	\notag\\
	&\eta^{11}\bigl(\Lambda_{m\mu},{\Gamma_m}^{\nu'}\bigr)=-\eta^{11}(\Lambda_{m\mu'},{\Gamma_m}^{\nu}) = \delta_{\mu\nu}.
	\label{BFosp2}
\end{align}

These bilinear forms are nondegenerate. Their inverses $\eta_{00}$ and $\eta_{11}$ are given by
\begin{align*}
	&\eta_{00}(T_{ij},T_{pq})= \eta_{00}({U_i}^j, {U_p}^q) = \eta_{11}(T_{ij},{U_p}^q) = \delta_{iq} \delta_{jp},
	\notag\\
	&\eta_{00}(T_{ij'}, T_{p'q})= \eta_{00}({U_i}^{j'}, {U_{p'}}^{q}) = \eta_{11}(T_{ij'},{U_{p'}}^q) = \eta_{11}(T_{i'j},{U_p}^{q'}) = -\delta_{ip} \delta_{jq},
	\notag \\
&	\eta_{00}(T_{im}, T_{j'm})= \eta_{00}({U_i}^m,{U_{j'}}^m) = \eta_{11}(T_{im},{U_{j'}}^m) = \eta_{11}(T_{i'm}, {U_j}^m) = -\delta_{ij},
	\notag \\
&	\eta_{00}(T_{\mu\nu}, T_{\rho\sigma})= \eta_{00}({U_{\mu}}^{\nu},{U_{\rho}}^{\sigma}) = -\eta_{11}(T_{\mu\nu},{U_{\rho}}^{\sigma})= -\delta_{\mu\sigma} \delta_{\nu\rho},
	\notag\\
&	\eta_{00}(T_{\mu\nu'}, T_{\rho'\sigma})= \eta_{00}({U_{\mu}}^{\nu'},{U_{\rho'}}^{\sigma}) = -\eta_{11}(T_{\mu\nu'},{U_{\rho'}}^{\sigma}) =
	-\eta_{11}(T_{\mu'\nu},{U_{\rho}}^{\sigma'})
	\notag\\
	&\hphantom{\eta_{00}(T_{\mu\nu'}, T_{\rho'\sigma})} {} =
	\begin{cases}
		-\delta_{\mu\rho} \delta_{\nu\sigma}, & \mu < \nu, \ \rho < \sigma,
		\\[5pt]
		-2 \delta_{\mu\rho}, & \mu = \nu, \ \rho=\sigma,
	\end{cases}
\end{align*}
and
\begin{align*}
	&\eta_{00}(\Lambda_{i\mu}, \Lambda_{j'\nu'})= -\eta_{00}(\Lambda_{i\mu'}, \Lambda_{j'\nu}) = -\eta_{00}\bigl({\Gamma_i}^{\mu}, {\Gamma_{j'}}^{\nu'}\bigr) = \eta_{00}({\Gamma_i}^{\mu'}, {\Gamma_{j'}}^{\nu}) = -\delta_{ij}\delta_{\mu\nu},
	\notag\\
&	\eta_{00}(\Lambda_{m\mu}, \Lambda_{m\nu'})= -\eta_{00}\bigl({\Gamma_m}^{\mu}, {\Gamma_m}^{\nu'}\bigr) = -\delta_{\mu\nu},
	\notag\\
&	\eta_{11}\bigl(\Lambda_{i\mu},{\Gamma_{j'}}^{\nu'}\bigr)= -\eta_{11}(\Lambda_{i\mu'},{\Gamma_{j'}}^{\nu}) = \eta_{11}\bigl(\Lambda_{i'\mu},{\Gamma_{j}}^{\nu'}\bigr) = -\eta_{11}(\Lambda_{i'\mu'},{\Gamma_j}^{\nu}) = \delta_{ij} \delta_{\mu\nu},
	\notag\\
&	\eta_{11}\bigl(\Lambda_{m\mu},{\Gamma_m}^{\nu'}\bigr)=-\eta_{11}(\Lambda_{m\mu'},{\Gamma_m}^{\nu}) = \delta_{\mu\nu}.
\end{align*}
Here again, additional nonvanishing components obtained from the symmetry properties are not listed.

We are now able to write down explicit formulas for the Casimir elements of \Zosp\ for $m = 2\ell+1$
\begin{align*}
	C_{00}={}& \sum_{i,j} \big( T_{ij}T_{ji} + {U_{i}}^{j}{U_{j}}^{i} \big) - \sum_{i<j} \big( \{T_{ij'},T_{i'j}\} + \{{U_{i}}^{j'},{U_{i'}}^{j}\} \big) \nn\\
	&-\sum_i \big( \{T_{im},T_{i'm}\} + \{{U_{i}}^m,{U_{i'}}^m\} \big)
	 - \sum_{\mu,\nu} \big( T_{\mu\nu}T_{\nu\mu} + {U_{\mu}}^{\nu}{U_{\nu}}^{\mu} \big) \nn\\
	&- 2\sum_\mu\big( \{T_{\mu\mu'},T_{\mu'\mu}\} + \{{U_{\mu}}^{\mu'},{U_{\mu'}}^{\mu}\} \big) - \sum_{\mu<\nu}\big( \{T_{\mu\nu'},T_{\mu'\nu}\} + \{{U_{\mu}}^{\nu'},{U_{\mu'}}^{\nu}\} \big) \nn\\
	&+ \sum_{i,\mu} \big( [\Lambda_{i'\mu'},\Lambda_{i\mu}] + [\Lambda_{i\mu'},\Lambda_{i'\mu}] + [{\Gamma_{i}}^{\mu},{\Gamma_{i'}}^{\mu'}] + [{\Gamma_{i'}}^{\mu},{\Gamma_{i}}^{\mu'}] \big) \nn\\
	&+ \sum_{\mu}\big( [\Lambda_{m\mu'},\Lambda_{m\mu}] + [{\Gamma_m}^{\mu},{\Gamma_m}^{\mu'}] \big),
\\
	C_{11}={}& \sum_{i,j} \{T_{ij},{U_j}^{i}\} - \sum_{i<j} \big( \{T_{ij'},{U_{i'}}^{j}\} + \{T_{i'j},{U_{i}}^{j'}\} \big) - \sum_{i}\big( \{T_{im},{U_{i'}}^m\} + \{T_{i'm},{U_{i}}^m\} \big) \nn\\
	&+ \sum_{\mu,\nu} \{T_{\mu\nu},{U_{\nu}}^{\mu}\} + 2 \sum_{\mu} \big( \{T_{\mu\mu'},{U_{\mu'}}^{\mu}\} + \{T_{\mu'\mu},{U_{\mu}}^{\mu'}\} \big) \nn\\
	&+ \sum_{\mu<\nu} \big( \{T_{\mu\nu'},{U_{\mu'}}^{\nu}\} + \{T_{\mu'\nu},{U_{\mu}}^{\nu'}\} \big) \nn\\
	&+ \sum_{i,\mu} \big( \{\Lambda_{i\mu},{\Gamma_{i'}}^{\mu'}\} - \{\Lambda_{i\mu'},{\Gamma_{i'}}^{\mu}\} + \{\Lambda_{i'\mu},{\Gamma_{i}}^{\mu'}\} - \{\Lambda_{i'\mu'},{\Gamma_{i}}^{\mu}\}\big) \nn\\
	&+ \sum_\mu \big( \{\Lambda_{m\mu},{\Gamma_{m}}^{\mu'}\} - \{\Lambda_{m\mu'},{\Gamma_{m}}^{\mu}\} \big).
\end{align*}
The corresponding Casimir elements for $m = 2\ell$ are obtained by eliminating the basis elements with index $m$.
Proposition~\ref{PROP:comm} also shows that the loop algebra of \Zosp\ admits two central extensions: one of degree $00$ and the other of degree~$11$.
Using the bilinear forms~\eqref{BFosp1} and~\eqref{BFosp2}, it is straightforward to write down all the commutation relations of the affine extension of \Zosp.
We therefore omit them here and simply note that all the relations for $\Z2$-$\osp(1|2)$ can be found in~\cite{AiSe}.

\begin{thebibliography}{99}
\footnotesize\itemsep=0pt

\bibitem{Aizawa_2020}
Aizawa N., Amakawa K., Doi S., {$\mathcal N$}-extension of double-graded
 supersymmetric and superconformal quantum mechanics,
 \href{https://doi.org/10.1088/1751-8121/ab661c}{\textit{J.~Phys.~A}}
 \textbf{53} (2020), 065205, 14~pages,
 \href{http://arxiv.org/abs/1905.06548}{arXiv:1905.06548}.

\bibitem{AiAmaDoiZn}
Aizawa N., Amakawa K., Doi S., {$\mathbb{Z}_2^n$}-graded extensions of
 supersymmetric quantum mechanics via {C}lifford algebras,
 \href{https://doi.org/10.1063/1.5144325}{\textit{J.~Math. Phys.}} \textbf{61}
 (2020), 052105, 13~pages,
 \href{http://arxiv.org/abs/1912.11195}{arXiv:1912.11195}.

\bibitem{AiFujiIto}
Aizawa N., Fujii I., Ito R., An integrable hierarchy associated with loop
 extension of {$\mathbb Z_2^2$}-graded {$\mathfrak{osp}(1|2)$},
 \href{https://doi.org/10.1088/1751-8121/ae6d53}{\textit{J.~Phys.~A}}
 \textbf{59} (2026), 215206, 15~pages,
 \href{http://arxiv.org/abs/2512.14108}{arXiv:2512.14108}.

\bibitem{AFITTsuper}
Aizawa N., Fujii I., Ito R., Tanaka T., Toppan F., Integrable {$\mathbb
 Z_2^2$}-graded super-{L}iouville equation and induced {$\mathbb
 Z_2^2$}-graded super-{V}irasoro algebra,
 \href{https://doi.org/10.1088/1751-8121/ae73c0}{\textit{J.~Phys.~A}}
 \textbf{59} (2026), 235202, 20~pages,
 \href{http://arxiv.org/abs/2512.17449}{arXiv:2512.17449}.

\bibitem{AizawaConfType}
Aizawa N., Isaac P.S., Segar J.,
 {$\mathbb{Z}_2\times\mathbb{Z}_2$}-generalizations of infinite-dimensional
 {L}ie superalgebra of conformal type with complete classification of central
 extensions, \href{https://doi.org/10.1016/S0034-4877(20)30041-0}{\textit{Rep.
 Math. Phys.}} \textbf{85} (2020), 351--373,
 \href{http://arxiv.org/abs/1902.05741}{arXiv:1902.05741}.

\bibitem{AiItoIntegral}
Aizawa N., Ito R., Integration on minimal {$\mathbb Z^2_2$}-superspace and
 emergence of space,
 \href{https://doi.org/10.1088/1751-8121/ad076e}{\textit{J.~Phys.~A}}
 \textbf{56} (2023), 485201, 11~pages,
 \href{http://arxiv.org/abs/2305.07836}{arXiv:2305.07836}.

\bibitem{AIKTTslint}
Aizawa N., Ito R., Kuznetsova Z., Tanaka T., Toppan F., Integrable {$\mathbb
 Z^2_2$}-graded extensions of the {L}iouville and {S}inh-{G}ordon theories,
 \href{https://doi.org/10.1088/1751-8121/adaab3}{\textit{J.~Phys.~A}}
 \textbf{58} (2025), 055201, 23~pages,
 \href{http://arxiv.org/abs/2406.13503}{arXiv:2406.13503}.

\bibitem{AiItoTa}
Aizawa N., Ito R., Tanaka T., {$\mathbb Z^2_2$}-graded supersymmetry via
 superfield on minimal {$\mathbb Z^2_2$}-superspace,
 \href{https://doi.org/10.1088/1751-8121/ad811a}{\textit{J.~Phys.~A}}
 \textbf{57} (2024), 435201, 19~pages,
 \href{http://arxiv.org/abs/2308.16860}{arXiv:2308.16860}.

\bibitem{AiKimura}
Aizawa N., Kimura D., Universal weight systems from a minimal {$\mathbb
 Z_2^2$}-graded {L}ie algebra,
 \href{https://doi.org/10.1142/S0218216525500452}{\textit{J.~Knot Theory
 Ramifications}} \textbf{34} (2025), 2550045, 30~pages,
 \href{http://arxiv.org/abs/2410.05845}{arXiv:2410.05845}.

\bibitem{LLEprogress}
Aizawa N., Kuznetsova Z., Tanaka H., Toppan F.,
 {$\mathbb{Z}_2\times\mathbb{Z}_2$}-graded {L}ie symmetries of the
 {L}\'{e}vy-{L}eblond equations,
 \href{https://doi.org/10.1093/ptep/ptw176}{\textit{PTEP. Prog. Theor. Exp.
 Phys.}} \textbf{2016} (2016), 123A01, 26~pages,
 \href{http://arxiv.org/abs/1609.08224}{arXiv:1609.08224}.

\bibitem{LLEproceed}
Aizawa N., Kuznetsova Z., Tanaka H., Toppan F., Generalized supersymmetry and
 the {L\'e}vy-{L}eblond equation, in Physical and Mathematical Aspects of
 Symmetries, \href{https://doi.org/10.1007/978-3-319-69164-0_11}{Springer},
 Cham, 2017, 79--84,
 \href{http://arxiv.org/abs/1609.08760}{arXiv:1609.08760}.

\bibitem{AiSe}
Aizawa N., Segar J., Affine extensions of {$\mathbb Z_2^2$}-graded {$\mathfrak{osp}(1|2)$}
 and {V}irasoro algebra,
 \href{https://doi.org/10.1142/S0219887825400523}{\textit{Int.~J. Geom.
 Methods Mod. Phys.}} \textbf{23} (2026), 2540052, 22~pages,
 \href{http://arxiv.org/abs/2409.07938}{arXiv:2409.07938}.

\bibitem{huerta2025particle}
Alderete C., Green A.M., Nguyen N.H., Zhu Y., Linke N.M., Rodr{\'\i}guez-Lara
 B.M., Para-particle oscillator simulations on a trapped-ion quantum computer,
 \href{https://doi.org/10.1063/5.0276426}{\textit{J.~Appl. Phys.}}
 \textbf{138} (2025), 054401, 8~pages,
 \href{http://arxiv.org/abs/2207.02430}{arXiv:2207.02430}.

\bibitem{alderete2025experimental}
Alderete C., Green A.M., Nguyen N.H., Zhu Y., Rodr{\'\i}guez-Lara B.M., Linke
 N.M., Experimental realization of para-particle oscillators,
 \href{https://doi.org/10.1038/s41598-025-25271-2}{\textit{Sci. Rep.}}
 \textbf{15} (2025), 41498, 11~pages,
 \href{http://arxiv.org/abs/2108.05471}{arXiv:2108.05471}.

\bibitem{AlHoSe}
Alshammari F.S., Hoque M.F., Segar J., The $\mathfrak{osp}(1|2)$ $\mathbb{Z}_2
 \times \mathbb{Z}_2$ graded algebra and its irreducible representations,
 \href{http://arxiv.org/abs/2306.12381}{arXiv:2306.12381}.

\bibitem{AmaAi}
Amakawa K., Aizawa N., A classification of lowest weight irreducible modules
 over {$\mathbb{Z}_2^2$}-graded extension of {$\mathfrak{osp}(1|2)$},
 \href{https://doi.org/10.1063/5.0037493}{\textit{J.~Math. Phys.}} \textbf{62}
 (2021), 043502, 18~pages,
 \href{http://arxiv.org/abs/2011.03714}{arXiv:2011.03714}.

\bibitem{BBTbook}
Babelon O., Bernard D., Talon M., Introduction to classical integrable systems,
 \textit{Cambridge Monogr. Math. Phys.},
 \href{https://doi.org/10.1017/CBO9780511535024}{Cambridge University Press},
 Cambridge, 2003.

\bibitem{BALBINO}
Balbino M.M., de~Freitas I.P., Rana R.G., Toppan F., Inequivalent
 {$\mathbb{Z}_2^n$}-graded brackets, {$n$}-bit parastatistics and statistical
 transmutations of supersymmetric quantum mechanics,
 \href{https://doi.org/10.1016/j.nuclphysb.2024.116729}{\textit{Nuclear
 Phys.~B}} \textbf{1009} (2024), 116729, 38~pages,
 \href{http://arxiv.org/abs/2309.00965}{arXiv:2309.00965}.

\bibitem{BruceSUSY}
Bruce A.J., On a $\mathbb{Z}_2^n$-graded version of supersymmetry,
 \href{https://doi.org/10.3390/sym11010116}{\textit{Symmetry}} \textbf{11}
 (2019), 116, 20~pages,
 \href{http://arxiv.org/abs/1812.02943}{arXiv:1812.02943}.

\bibitem{BruSigma}
Bruce A.J., {$\mathbb Z_2 \times \mathbb Z_2$}-graded supersymmetry: 2-d sigma
 models, \href{https://doi.org/10.1088/1751-8121/abb47f}{\textit{J.~Phys.~A}}
 \textbf{53} (2020), 455201, 25~pages,
 \href{http://arxiv.org/abs/2006.08169}{arXiv:2006.08169}.

\bibitem{BruSG}
Bruce A.J., Is the {$\mathbb Z_2\times\mathbb Z_2$}-graded sine-{G}ordon
 equation integrable?,
 \href{https://doi.org/10.1016/j.nuclphysb.2021.115514}{\textit{Nuclear
 Phys.~B}} \textbf{971} (2021), 115514, 10~pages,
 \href{http://arxiv.org/abs/2106.06372}{arXiv:2106.06372}.

\bibitem{BruDup}
Bruce A.J., Duplij S., Double-graded supersymmetric quantum mechanics,
 \href{https://doi.org/10.1063/1.5118302}{\textit{J.~Math. Phys.}} \textbf{61}
 (2020), 063503, 13~pages,
 \href{http://arxiv.org/abs/1904.06975}{arXiv:1904.06975}.

\bibitem{BruIb}
Bruce A.J., Ibarguengoytia E., Poncin N., The {S}chwarz--{V}oronov embedding of
 {$\mathbb Z_2^n$}-manifolds,
 \href{https://doi.org/10.3842/SIGMA.2020.002}{\textit{SIGMA}} \textbf{16}
 (2020), 002, 47~pages,
 \href{http://arxiv.org/abs/1906.09834}{arXiv:1906.09834}.

\bibitem{BruPon}
Bruce A.J., Poncin N., Functional analytic issues in {$\mathbb
 Z_2^n$}-geometry,
 \href{https://doi.org/10.33044/revuma.v60n2a21}{\textit{Rev. Un. Mat.
 Argentina}} \textbf{60} (2019), 611--636,
 \href{http://arxiv.org/abs/1807.11739}{arXiv:1807.11739}.

\bibitem{bruce2019products}
Bruce A.J., Poncin N., Products in the category of {$\mathbb Z^n_2$}-manifolds,
 \href{https://doi.org/10.1080/14029251.2019.1613051}{\textit{J.~Nonlinear
 Math. Phys.}} \textbf{26} (2019), 420--453,
 \href{http://arxiv.org/abs/1807.11740}{arXiv:1807.11740}.

\bibitem{Covolo}
Covolo T., Cohomological approach to the graded {B}erezinian,
 \href{https://doi.org/10.4171/JNCG/200}{\textit{J.~Noncommut. Geom.}}
 \textbf{9} (2015), 543--565,
 \href{http://arxiv.org/abs/1207.2962}{arXiv:1207.2962}.

\bibitem{CoGrPo}
Covolo T., Grabowski J., Poncin N., The category of
 {$\mathbb{Z}_2^n$}-supermanifolds,
 \href{https://doi.org/10.1063/1.4955416}{\textit{J.~Math. Phys.}} \textbf{57}
 (2016), 073503, 16~pages,
 \href{http://arxiv.org/abs/1602.03312}{arXiv:1602.03312}.

\bibitem{CoGrPo2016}
Covolo T., Grabowski J., Poncin N., Splitting theorem for
 $\mathbb{Z}_2^n$-supermanifolds,
 \href{https://doi.org/10.1016/j.geomphys.2016.09.006}{\textit{J.~Geom.
 Phys.}} \textbf{110} (2016), 393--401,
 \href{http://arxiv.org/abs/1602.03671}{arXiv:1602.03671}.

\bibitem{CoKwoPo}
Covolo T., Kwok S., Poncin N., Local forms of morphisms of colored
 supermanifolds,
 \href{https://doi.org/10.1016/j.geomphys.2021.104302}{\textit{J.~Geom.
 Phys.}} \textbf{168} (2021), 104302, 21~pages,
 \href{http://arxiv.org/abs/2010.10026}{arXiv:2010.10026}.

\bibitem{CoGrPo2012}
Covolo T., Ovsienko V., Poncin N., Higher trace and {B}erezinian of matrices
 over a {C}lifford algebra,
 \href{https://doi.org/10.1016/j.geomphys.2012.07.004}{\textit{J.~Geom.
 Phys.}} \textbf{62} (2012), 2294--2319,
 \href{http://arxiv.org/abs/1109.5877}{arXiv:1109.5877}.

\bibitem{GrJav}
Green H.S., Jarvis P.D., Casimir invariants, characteristic identities, and
 {Y}oung diagrams for color algebras and superalgebras,
 \href{https://doi.org/10.1063/1.525911}{\textit{J.~Math. Phys.}} \textbf{24}
 (1983), 1681--1687.

\bibitem{IsRyan}
Isaac P.S., Ryan M., The {D}irac oscillator, generalised parastatistics and
 colour {L}ie superalgebras,
 \href{https://doi.org/10.1016/j.nuclphysb.2025.117154}{\textit{Nuclear
 Phys.~B}} \textbf{1020} (2025), 117154, 15~pages,
 \href{http://arxiv.org/abs/2507.14780}{arXiv:2507.14780}.

\bibitem{IsStVdJ}
Isaac P.S., Stoilova N.I., Van~der Jeugt J., The
 {$\mathbb{Z}_2\times\mathbb{Z}_2$}-graded general linear {L}ie superalgebra,
 \href{https://doi.org/10.1063/1.5138597}{\textit{J.~Math. Phys.}} \textbf{61}
 (2020), 011702, 7~pages,
 \href{http://arxiv.org/abs/1912.08636}{arXiv:1912.08636}.

\bibitem{ItoNago}
Ito R., Nago A., Novel possible symmetries of {$S$}-matrix generated by
 {$\mathbb Z_2^n$}-graded {L}ie superalgebras,
 \href{https://doi.org/10.1016/j.nuclphysb.2025.116877}{\textit{Nuclear
 Phys.~B}} \textbf{1014} (2025), 116877, 10~pages,
 \href{http://arxiv.org/abs/2501.07311}{arXiv:2501.07311}.

\bibitem{JaYangWyb}
Jarvis P.D., Yang M., Wybourne B.G., Generalized quasispin for supergroups,
 \href{https://doi.org/10.1063/1.527566}{\textit{J.~Math. Phys.}} \textbf{28}
 (1987), 1192--1197.

\bibitem{KuToclassification}
Kuznetsova Z., Toppan F., Classification of minimal {$\mathbb Z_2\times\mathbb
 Z_2$}-graded {L}ie (super)algebras and some applications,
 \href{https://doi.org/10.1063/5.0050200}{\textit{J.~Math. Phys.}} \textbf{62}
 (2021), 063512, 22~pages,
 \href{http://arxiv.org/abs/2103.04385}{arXiv:2103.04385}.

\bibitem{LU20231}
Lu R., Tan Y., Construction of color {L}ie algebras from homomorphisms of
 modules of {L}ie algebras,
 \href{https://doi.org/10.1016/j.jalgebra.2022.12.028}{\textit{J.~Algebra}}
 \textbf{620} (2023), 1--49.

\bibitem{MoVar}
Mohammadi M., Varsaie S., On the construction of {$\mathbb
 Z^n_2$}-grassmannians as homogeneous {$\mathbb Z^n_2$}-spaces,
 \href{https://doi.org/10.3934/era.2022012}{\textit{Electron. Res. Arch.}}
 \textbf{30} (2022), 221--241,
 \href{http://arxiv.org/abs/1901.07393}{arXiv:1901.07393}.

\bibitem{SilvCoho}
Piontkovski D., Silvestrov S.D., Cohomology of 3-dimensional color {L}ie
 algebras,
 \href{https://doi.org/10.1016/j.jalgebra.2006.11.008}{\textit{J.~Algebra}}
 \textbf{316} (2007), 499--513,
 \href{http://arxiv.org/abs/math.KT/0508573}{arXiv:math.KT/0508573}.

\bibitem{Poncin}
Poncin N., Towards integration on colored supermanifolds, in Geometry of Jets
 and Fields, \textit{Banach Center Publ.}, Vol. 110, Polish Academy of
 Sciences, Institute of Mathematics, Warsaw, 2016, 201--217.

\bibitem{PonSch}
Poncin N., Schouten S., The geometry of supersymmetry: {A} concise
 introduction, \textit{Grad.~J. Math.} \textbf{8} (2023), 1--57,
 \href{http://arxiv.org/abs/2207.12974}{arXiv:2207.12974}.

\bibitem{Ree}
Ree R., Generalized {L}ie elements,
 \href{https://doi.org/10.4153/CJM-1960-044-x}{\textit{Canadian~J. Math.}}
 \textbf{12} (1960), 493--502.

\bibitem{rw1}
Rittenberg V., Wyler D., Generalized superalgebras,
 \href{https://doi.org/10.1016/0550-3213(78)90186-4}{\textit{Nuclear Phys.~B}}
 \textbf{139} (1978), 189--202.

\bibitem{rw2}
Rittenberg V., Wyler D., Sequences of {$Z_{2}\oplus Z_{2}$} graded {L}ie
 algebras and superalgebras,
 \href{https://doi.org/10.1063/1.523552}{\textit{J.~Math. Phys.}} \textbf{19}
 (1978), 2193--2200.

\bibitem{RyanRefining}
Ryan M., Refining the grading of irreducible {L}ie colour algebra
 representations, \href{http://arxiv.org/abs/2403.02855}{arXiv:2403.02855}.

\bibitem{RyanLL}
Ryan M., Graded colour {L}ie superalgebras for solving {L}\'{e}vy-{L}eblond
 equations,
 \href{https://doi.org/10.1088/1751-8121/ad978a}{\textit{J.~Phys.~A}}
 \textbf{58} (2025), 015204, 24~pages,
 \href{http://arxiv.org/abs/2407.19723}{arXiv:2407.19723}.

\bibitem{scheu}
Scheunert M., Generalized {L}ie algebras,
 \href{https://doi.org/10.1063/1.524113}{\textit{J.~Math. Phys.}} \textbf{20}
 (1979), 712--720.

\bibitem{scheuCasi}
Scheunert M., Casimir elements of {$\varepsilon $} {L}ie algebras,
 \href{https://doi.org/10.1063/1.525642}{\textit{J.~Math. Phys.}} \textbf{24}
 (1983), 2671--2680.

\bibitem{scheuTen}
Scheunert M., Graded tensor calculus,
 \href{https://doi.org/10.1063/1.525641}{\textit{J.~Math. Phys.}} \textbf{24}
 (1983), 2658--2670.

\bibitem{SchZhang}
Scheunert M., Zhang R.B., Cohomology of {L}ie superalgebras and their
 generalizations, \href{https://doi.org/10.1063/1.532508}{\textit{J.~Math.
 Phys.}} \textbf{39} (1998), 5024--5061,
 \href{http://arxiv.org/abs/q-alg/9701037}{arXiv:q-alg/9701037}.

\bibitem{Sergeev1983}
Sergeev A.N., The centre of enveloping algebra for {L}ie superalgebra
 {$Q(n,\mathbb{C})$}, \href{https://doi.org/10.1007/BF00400431}{\textit{Lett.
 Math. Phys.}} \textbf{7} (1983), 177--179.

\bibitem{SVdJ2018}
Stoilova N.I., Van~der Jeugt J., The {$\mathbb Z_2\times\mathbb Z_2$}-graded
 {L}ie superalgebra {$\mathfrak{pso}(2m+1|2n)$} and new parastatistics
 representations,
 \href{https://doi.org/10.1088/1751-8121/aaae9a}{\textit{J.~Phys.~A}}
 \textbf{51} (2018), 135201, 17~pages,
 \href{http://arxiv.org/abs/1711.02136}{arXiv:1711.02136}.

\bibitem{Stclassical}
Stoilova N.I., Van~der Jeugt J., On classical {$\mathbb Z_2\times\mathbb
 Z_2$}-graded {L}ie algebras,
 \href{https://doi.org/10.1063/5.0149175}{\textit{J.~Math. Phys.}} \textbf{64}
 (2023), 061702, 8~pages,
 \href{http://arxiv.org/abs/2305.18604}{arXiv:2305.18604}.

\bibitem{SVdJ2024}
Stoilova N.I., Van~der Jeugt J., Orthosymplectic {$\mathbb Z_2\times\mathbb
 Z_2$}-graded {L}ie superalgebras and parastatistics,
 \href{https://doi.org/10.1088/1751-8121/ad2726}{\textit{J.~Phys.~A}}
 \textbf{57} (2024), 095202, 13~pages,
 \href{http://arxiv.org/abs/2402.11952}{arXiv:2402.11952}.

\bibitem{StG2}
Stoilova N.I., Van~der Jeugt J., {$\mathbb Z^3_2$}-grading of the {L}ie algebra
 {$G_2$} and related color algebras,
 \href{https://doi.org/10.1088/1751-8121/adfea7}{\textit{J.~Phys.~A}}
 \textbf{58} (2025), 365201, 17~pages,
 \href{http://arxiv.org/abs/2505.04378}{arXiv:2505.04378}.

\bibitem{tolstoyOnce}
Tolstoy V.N., Once more on parastatistics,
 \href{https://doi.org/10.1134/S1547477114070449}{\textit{Phys. Part. Nucl.
 Lett.}} \textbf{11} (2014), 933--937,
 \href{http://arxiv.org/abs/1610.01628}{arXiv:1610.01628}.

\bibitem{tolstoy2014}
Tolstoy V.N., Super--de {S}itter and alternative super-{P}oincar\'{e}
 symmetries, in Lie Theory and Its Applications in Physics, \textit{Springer
 Proc. Math. Stat.}, Vol. 111,
 \href{https://doi.org/10.1007/978-4-431-55285-7_26}{Springer}, Tokyo, 2014,
 357--367, \href{http://arxiv.org/abs/1610.01566}{arXiv:1610.01566}.

\bibitem{toppan2021inequivalent}
Toppan F., Inequivalent quantizations from gradings and
 {$\mathbb{Z}_2\times\mathbb{Z}_2$} parabosons,
 \href{https://doi.org/10.1088/1751-8121/ac17a5}{\textit{J.~Phys.~A}}
 \textbf{54} (2021), 355202, 21~pages,
 \href{http://arxiv.org/abs/2104.09692}{arXiv:2104.09692}.

\bibitem{toppanMulti}
Toppan F., {$\mathbb Z_2 \times \mathbb Z_2$}-graded parastatistics in
 multiparticle quantum {H}amiltonians,
 \href{https://doi.org/10.1088/1751-8121/abe2f2}{\textit{J.~Phys.~A}}
 \textbf{54} (2021), 115203, 35~pages,
 \href{http://arxiv.org/abs/2008.11554}{arXiv:2008.11554}.

\bibitem{toppan2022first}
Toppan F., First quantization of braided {M}ajorana fermions,
 \href{https://doi.org/10.1016/j.nuclphysb.2022.115834}{\textit{Nuclear
 Phys.~B}} \textbf{980} (2022), 115834, 14~pages,
 \href{http://arxiv.org/abs/2203.01776}{arXiv:2203.01776}.

\bibitem{toppan2024braid}
Toppan F., On braid statistics versus parastatistics,
 \href{https://doi.org/10.1088/1742-6596/2912/1/012011}{\textit{J.~Phys. Conf.
 Ser.}} \textbf{2912} (2024), 012011, 11~pages,
 \href{http://arxiv.org/abs/2411.14261}{arXiv:2411.14261}.

\bibitem{toppan2024detectability}
Toppan F., On the detectability of paraparticles beyond bosons and fermions,
 \href{https://doi.org/10.1142/S0219887825400420}{\textit{Int.~J. Geom.
 Methods Mod. Phys.}} \textbf{23} (2026), 2540042, 21~pages,
 \href{http://arxiv.org/abs/2411.18313}{arXiv:2411.18313}.

\bibitem{Vasiliev}
Vasiliev M.A., de {S}itter supergravity with positive cosmological constant and
 generalised {L}ie superalgebras,
 \href{https://doi.org/10.1088/0264-9381/2/5/007}{\textit{Classical Quantum
 Gravity}} \textbf{2} (1985), 645--659.

\bibitem{wang2025particle}
Wang Z., Hazzard K.R.A., Particle exchange statistics beyond fermions and
 bosons, \href{https://doi.org/10.1038/s41586-024-08262-7}{\textit{Nature}}
 \textbf{637} (2025), 314--318,
 \href{http://arxiv.org/abs/2308.05203}{arXiv:2308.05203}.

\bibitem{WTIQ}
Wills~Toro L.A., {$(I,q)$}-graded {L}ie algebraic extensions of the
 {P}oincar\'{e} algebra, constraints on {$I$} and {$q$},
 \href{https://doi.org/10.1063/1.531102}{\textit{J.~Math. Phys.}} \textbf{36}
 (1995), 2085--2112.

\bibitem{WTz4}
Wills-Toro L.A., Trefoil symmetries. {I}. {C}lover extensions beyond
 {C}oleman--{M}andula theorem,
 \href{https://doi.org/10.1063/1.1383561}{\textit{J.~Math. Phys.}} \textbf{42}
 (2001), 3915--3934.

\bibitem{zhang2025}
Zhang R.B., Invariants and representations of the {$\Gamma$}-graded general
 linear {L}ie {$\omega$}-algebras,
 \href{https://doi.org/10.1016/j.exmath.2026.125778}{\textit{Expo. Math.}}
 \textbf{44} (2026), 125778, 92~pages,
 \href{http://arxiv.org/abs/2509.21795}{arXiv:2509.21795}.

\bibitem{zhang2026}
Zhang R.B., Quantum groups of {L}ie colour algebras fulfilling {C}artan--{W}eyl
 paradigm, \href{http://arxiv.org/abs/2606.02202}{arXiv:2606.02202}.

\end{thebibliography}
\end{document}
